
% !TEX program = pdflatex
\documentclass[12pt,a4paper,oneside]{report}
\usepackage[utf8]{inputenc}
\usepackage[T1]{fontenc}
\usepackage{lmodern}
\usepackage{CJKutf8}
\usepackage{geometry}
\geometry{left=2.8cm,right=2.8cm,top=2.7cm,bottom=2.7cm}
\usepackage{graphicx}
\usepackage{float}
\usepackage{booktabs}
\usepackage{array}
\usepackage{longtable}
\usepackage{amsmath,amssymb}
\usepackage{setspace}
\usepackage{hyperref}
\usepackage{caption}
\usepackage{subcaption}
\usepackage{xcolor}
\usepackage{enumitem}
\usepackage{multirow}
\usepackage{url}
\usepackage{siunitx}
\usepackage{listings}
\usepackage{fancyhdr}
\usepackage{titlesec}
\usepackage{microtype}

\graphicspath{{figures/}}
\hypersetup{colorlinks=true,linkcolor=black,citecolor=blue,urlcolor=blue}
\onehalfspacing
\sloppy
\emergencystretch=3em
\setlength{\parindent}{2em}
\setlength{\parskip}{0.25em}
\setlength{\headheight}{15pt}
\captionsetup{font=small,labelfont=bf}
\renewcommand{\bibname}{References}
\pagestyle{fancy}
\fancyhf{}
\fancyhead[L]{Simulation of~\textit{Staphylococcus aureus}~biofilms}
\fancyhead[R]{\thepage}
\newcommand{\zhtext}[1]{\begin{CJK*}{UTF8}{gbsn}#1\end{CJK*}}
\newcommand{\safeincludegraphics}[2][]{%
  \IfFileExists{#2}{%
    \includegraphics[#1]{#2}%
  }{%
    \fbox{%
      \begin{minipage}[c][0.24\textheight][c]{0.82\linewidth}
      \centering
      Missing figure: \texttt{\detokenize{#2}}\\
      \small Upload this file into the Overleaf \texttt{figures/} folder.
      \end{minipage}%
    }%
  }%
}

\titleformat{\chapter}{\normalfont\bfseries\Huge}{Chapter \thechapter}{1em}{}
\titleformat{\section}{\normalfont\bfseries\Large}{\thesection}{0.8em}{}
\titleformat{\subsection}{\normalfont\bfseries\large}{\thesubsection}{0.8em}{}

\lstset{
  basicstyle=\ttfamily\footnotesize,
  breaklines=true,
  frame=single,
  backgroundcolor=\color{gray!6},
  columns=fullflexible
}

\begin{document}

\begin{titlepage}
\centering
\vspace*{0.5cm}
\begin{minipage}{0.42\textwidth}
\centering
\safeincludegraphics[width=0.55\textwidth]{figures/original_word_figure_2.png}
\end{minipage}
\begin{minipage}{0.42\textwidth}
\centering
\safeincludegraphics[width=0.40\textwidth]{figures/original_word_figure_3.png}
\end{minipage}

\vspace{1.2cm}
{\large Graduation Thesis for SWJTU student enrolled in the SWJTU-GSU collaborative dual degree program\par}
\vspace{1.3cm}
{\Huge\bfseries Simulation of \textit{Staphylococcus aureus} biofilms and data-driven mechanism analysis\par}
\vspace{0.4cm}
{\Large \zhtext{金黄色葡萄球菌生物膜的仿真与数据驱动机制分析}\par}
\vspace{1.4cm}
{\large In partial fulfill requirements of SWJTU bachelor’s degree\par}
\vspace{1.1cm}
\begin{tabular}{rl}
Name: & Haolin Lyu\\
Major: & Bioengineering\\
Grade: & 2022\\
Student ID at SWJTU: & 2022113670\\
Panther ID at GSU: & 002819989\\
Supervisor at SWJTU: & Chun Ye\\
Supervisor at GSU: & Yi Jiang\\
Date: & May 20th, 2026\\
\end{tabular}
\vfill
\end{titlepage}

\pagenumbering{roman}

\chapter*{Abstract}
\addcontentsline{toc}{chapter}{Abstract}
In recent years, with the increasing prominence of catheter-related infection issues, the formation mechanism and structural evolution laws of biofilms in complex fluid environments have gradually become research hotspots. Traditional studies have mostly focused on the growth process of biofilms at a single scale, while there is still a lack of systematic analysis of the coupling mechanism between cell behavior, fluid dynamics, and environmental diffusion processes. At the same time, how to convert multi-scale simulation results into data-driven models that can be used for prediction and analysis also requires further exploration.

This study, based on the context of the catheter environment, constructed a multi-scale biofilm simulation framework that combines cellular-level behaviors with the macroscopic fluid environment. By introducing the Cellular Potts Model (CPM) to describe the processes of cell growth and division, and combining the Navier--Stokes equations with diffusion models to simulate the fluid and nutrient/signal transmission processes, a systematic modeling of the formation and dynamic evolution of biofilms was achieved. On this basis, multi-dimensional feature information was extracted from the simulation outputs, a structured feature matrix was constructed, and a deep learning model was used to predict and analyze key structural indicators such as the thickness of the biofilm.

The results show that the constructed multi-layer perceptron model can capture the nonlinear relationship between the input features and the structure of the biological membrane to a certain extent. During the training process, the loss function of the model steadily decreased, and the performance on the validation set was consistent with that on the training set, indicating that the model has a certain generalization ability. At the same time, the prediction results show that the model can reflect the overall trend of the thickness of the biological membrane, but there is still a certain deviation in the local prediction accuracy, indicating that the current model still has room for improvement in terms of data scale and feature representation.

This study is the first to systematically integrate multi-scale biofilm simulation with data-driven modeling methods, proposing a unified research framework from mechanism modeling to predictive analysis. This method provides a new research approach for understanding the evolution mechanism of biofilm structures in fluid environments, and offers a theoretical basis for assessing the risk of catheter-related infections and formulating biofilm control strategies.

\noindent\textbf{Keywords:} Biofilm, multi-scale modeling, deep learning, fluid dynamics, feature engineering, catheter infection

\vspace{0.8cm}
The completed implementation further extends the original simulation-oriented thesis into a reproducible data platform. The final workflow links biofilm simulation, database storage, feature engineering, exploratory analysis, traditional machine learning, deep learning, and an interactive prediction interface. In this sense, the central contribution is not only a simulator, but a pipeline that converts cell-level simulation records into biofilm-level predictive evidence. This distinction is important: a single bacterium is not treated as an independent training sample. Instead, cell and cluster trajectories are aggregated by \texttt{simulation\_id} and timepoint, early-stage structural descriptors are extracted, and each complete simulation becomes one sample for final outcome prediction. This design makes the analysis closer to the biological question, because the targets of interest are biofilm-level properties such as thickness, roughness, surface coverage, EPS fraction, cluster count, and the time required for a stable biofilm-like state to emerge.

\tableofcontents
\clearpage

% Original Word document contained the placeholder text "空 页". In the LaTeX version it is represented as a blank page.
\thispagestyle{empty}\mbox{}\clearpage

\pagenumbering{arabic}

\chapter{Introduction}

\section{Background and Significance of Biofilm Research}
A biofilm is a complex three-dimensional structure system formed by microorganisms adhering to a solid surface and secreting extracellular polymeric substances (EPS) \cite{donlan2002}. Compared with free-floating microorganisms, the cells within a biofilm possess stronger environmental adaptability and resilience, especially in terms of antibiotic tolerance and immune evasion \cite{hallstoodley2004}. Therefore, biofilms have significant research value in environmental engineering, industrial systems, and clinical medicine.

\begin{figure}[H]
\centering
\safeincludegraphics[width=0.78\textwidth]{figures/original_word_figure_4.png}
\caption{Three-dimensional structure diagram of the biological membrane under microscopy. This image was preserved from the original Word thesis file and is used here as the introductory visual material.}
\label{fig:original-biofilm-micrograph}
\end{figure}

In recent years, with the widespread application of medical devices, catheter-associated infections (Catheter-Associated Infections) have gradually become an important issue in clinical practice \cite{costerton1999}. The formation of biofilms on the inner walls of catheters not only leads to the persistence of infections but also may cause pipe blockages and equipment failures. Therefore, studying the formation mechanism and evolution process of biofilms in fluid environments is of great significance for understanding their growth patterns and formulating effective control strategies.

\subsection{Catheter Environment and Biofilm Issues}
The internal environment of the conduit exhibits typical fluid characteristics. The velocity distribution, shear force, and material transfer processes within it all have significant impacts on the behavior of microorganisms. In actual systems, the fluid not only carries nutrients but also influences the attachment, detachment, and redistribution processes of bacteria.

On the surface of the catheter, bacteria initially attach to the surface through physical or chemical means. Subsequently, they multiply and secrete EPS to form micro-colony structures. Over time, the biofilm gradually develops into a mature structure and may undergo local detachment and reattachment under the influence of fluid. This dynamic process results in a highly uneven spatial distribution of the biofilm, and also increases the difficulty of control.

Therefore, it is necessary to conduct a systematic study on the formation and evolution process of the biomembrane within the duct from the perspective of ``fluid-cell behavior coupling''.

\begin{figure}[H]
\centering
\safeincludegraphics[width=0.88\textwidth]{figures/original_word_figure_5.png}
\caption{Diagram illustrating the process of biofilm formation in a catheter or pipe-like environment. The image was retained from the original thesis document.}
\label{fig:catheter-biofilm-process}
\end{figure}

\subsection{Mechanism of Biofilm Formation}
The formation of a biofilm typically involves three main stages: initial attachment, growth and formation, and dynamic evolution.

Firstly, during the initial attachment stage, bacteria attach to the inner wall of the tube through surface interactions, such as van der Waals forces and electrostatic forces. This process is influenced by both fluid shear force and surface properties.

Secondly, during the stage of biofilm formation, cells start to proliferate and secrete EPS, gradually forming micro-colonies with a certain structure. EPS not only provides mechanical support for the cells, but also regulates the transmission of nutrients and signaling molecules.

Finally, during the dynamic evolution stage, the biofilm may undergo local detachment under the influence of the fluid, while new bacteria continuously attach and form new structures. This ``growth--detachment--redistribution'' cycle enables the biofilm to maintain a dynamic equilibrium state.

\begin{figure}[H]
\centering
\safeincludegraphics[width=0.72\textwidth]{figures/original_word_figure_6.png}
\caption{Biofilm formation schematic preserved from the original Word thesis file. The schematic shows attachment, microcolony formation, matrix formation, maturation, and dispersion.}
\label{fig:biofilm-formation-schematic}
\end{figure}

\subsection{Impact of Fluid Environment on Biofilms}
The fluid environment is one of the important factors influencing the formation of biofilms. The fluid flow within the conduits can directly act on the cells through shear force, altering their attachment stability \cite{stoodley1994}. At the same time, the fluid also controls the transmission processes of nutrients, metabolic products, and signaling molecules.

Under low flow velocity conditions, bacteria are more likely to adhere and form stable structures; while under high flow velocity conditions, although the adhesion difficulty increases, the material transfer efficiency is higher, which may promote rapid local growth. Moreover, the non-uniformity of the flow field, such as vortices or stagnant zones, will form ``growth advantage areas'', thereby leading to spatial heterogeneity of the biofilm structure.

Therefore, incorporating fluid dynamics processes into the modeling of biological membranes is crucial for accurately describing their behavior.

\subsection{Overview of Multi-scale Modeling Methods}
Since the biomembrane system involves multiple processes such as cell behavior, substance transport, and fluid dynamics, its modeling usually requires the combination of multi-scale methods.

\subsection{Cellular Potts Model (CPM)}
The Cellular Potts Model is a lattice-based modeling method, which is often used to simulate the morphological changes, migration, and interactions of cells. In this study, CPM was employed to describe the growth, division, and formation of spatial organizational structures of bacteria.

\begin{figure}[H]
\centering
\safeincludegraphics[width=0.45\textwidth]{figures/original_word_figure_7.png}
\caption{A schematic diagram of cell behavior simulation based on the Cellular Potts Model, used to describe cell growth and interactions. This image was preserved from the original thesis document.}
\label{fig:cpm-original}
\end{figure}

\subsection{Navier--Stokes Fluid Model}
The Navier--Stokes equations are the fundamental equations for describing fluid motion and can be used to simulate the velocity distribution and pressure changes of fluids within conduits. In this study, the flow field was calculated using this model to analyze the impact of fluids on the structure of biological membranes.

\subsection{Diffusion and Convection Model}
In the biological membrane system, the transmission of nutrients and signaling molecules is usually determined by both diffusion and convection. By establishing a diffusion-convection model, the distribution of substances in space and their regulatory effects on cellular behavior can be described.

\subsection{Application of Deep Learning in Biofilm Research}
With the development of data-driven methods, deep learning has been widely applied in the field of biomedicine. Compared with traditional models, deep learning can extract nonlinear relationships from complex data, thereby improving the predictive ability.

In the field of biofilm research, simulation models can generate a large amount of spatiotemporal data, such as cell counts, spatial distribution, fluid parameters, etc. These data can be used as input for deep learning models to predict the growth trend or key indicators of biofilms.

Therefore, integrating physical modeling with deep learning is expected to enable a leap from ``understanding the mechanism'' to ``data prediction''.

\subsection{Research Questions and Research Objectives}
Based on the above background, this study aims to explore the interaction between the fluid environment within the catheter and the behavior of cells, and to analyze its impact on the formation and structural evolution of biological membranes. The specific research objectives include:

\begin{enumerate}[label=(\arabic*)]
\item Construct a biofilm growth model based on CPM to simulate the behavior of cells and the formation process of the structure;
\item Introduce the Navier--Stokes equation to simulate the fluid flow within the duct and its impact on the biofilm;
\item Establish a diffusion and convection model to describe the process of substance transfer;
\item Based on the simulation data, construct a feature matrix and use a deep learning model for predictive analysis;
\item Evaluate the model performance and analyze its potential application in actual biofilm research.
\end{enumerate}

Through the above research, this work aims to provide new theoretical support for understanding the formation mechanism of biological membranes in complex fluid environments, and to offer reference basis for future control strategies and engineering applications.

\section{Extension of the Original Research Logic}
The original introduction already defines the biological and modeling motivation. The implemented project extends that plan into a full data-driven research platform. The completed workflow can be summarized as follows: biofilm simulation produces raw spatiotemporal records; these records are exported into standard relational tables; cell-level and cluster-level observations are aggregated into biofilm-level descriptors; machine learning and deep learning models learn the mapping from early-stage conditions to final morphology; and an interactive application allows new parameter combinations or early partial simulation data to be evaluated without manually repeating every downstream analysis step.

It is worth noting that the computational objective is not simply to make a visually interesting simulation. A simulation alone can show how cells move and divide, but it does not automatically tell us which variables matter most, whether early differences predict later structure, or how reliable a prediction may be for a new parameter set. For this reason, the platform treats the simulator as a data-generating instrument. Each simulation run becomes a controlled experiment. The database then acts as the laboratory notebook, preserving the input parameters, process-level observations, flow-field descriptors, and final outcomes. This organization makes it possible to compare many virtual experiments in a systematic way.

The research question can therefore be refined into three connected questions. First, can a simplified but mechanistically interpretable simulation reproduce important qualitative features of biofilm development, such as wall attachment, cell proliferation, cluster formation, EPS-producing states, and shear-related detachment? Second, can the simulation output be transformed into a feature matrix that represents biofilm-level dynamics rather than isolated single-cell records? Third, can data-driven models use these features to predict final biofilm properties and reveal which early signals or physical parameters are most strongly associated with the resulting structure? These questions guided the design of the later chapters.


\section{Biological Context of \textit{Staphylococcus aureus} Biofilms}
\textit{Staphylococcus aureus} is a clinically important opportunistic pathogen. It is able to colonize skin, mucosal surfaces, wounds, and implanted medical devices. Once it forms a biofilm on a device surface, treatment becomes more difficult because the bacterial community is no longer exposed to the environment as isolated planktonic cells. Cells inside a biofilm may experience restricted antimicrobial penetration, altered metabolism, local nutrient gradients, and matrix-associated protection. This is one reason catheter-associated infections are often persistent and recurrent.

The biofilm lifestyle of \textit{S. aureus} is not a single event. It involves initial surface contact, adhesion strengthening, proliferation, cell-cell accumulation, EPS or matrix-associated development, maturation, and possible dispersal. Different strains and environmental conditions may emphasize different mechanisms. Some biofilms are strongly polysaccharide-dependent, while others may rely more on proteins or extracellular DNA. In this thesis, the simulation does not attempt to represent all molecular pathways separately. Instead, it abstracts them into state transitions and structural variables that can be measured and modeled.

This abstraction is necessary because a thesis-scale computational platform must balance biological realism and computational tractability. A fully molecular model would require many parameters that are difficult to estimate. A purely statistical model would ignore the mechanistic processes that make biofilm formation biologically meaningful. The chosen approach sits between these extremes. It represents cells as individual agents, includes fluid and signal effects, and then uses data-driven analysis to examine the generated patterns.

\section{From Visual Simulation to Quantitative Research Platform}
Many simulation projects stop at visual output. A movie or screenshot can show that a biofilm-like structure appears, but it cannot easily answer questions such as which parameter matters most, how early dynamics relate to final structure, or whether a new parameter combination will produce a thick biofilm. The current project was therefore developed as a quantitative platform.

The platform has three layers. The first layer is the simulation layer, where cell states, flow, signal diffusion, attachment, division, and EPS behavior are updated. The second layer is the data layer, where raw observations are stored in standardized tables. The third layer is the analysis and prediction layer, where features are engineered and models are trained. This layered design is important because it separates mechanism generation from statistical interpretation. If a figure looks unusual, the researcher can trace it back to raw database rows. If a model performs poorly, the researcher can check whether the issue comes from the feature matrix or the target distribution.

This design also makes the project easier to extend. A new simulation rule can be added without rewriting the database schema if the exported fields remain consistent. A new model can be tested without rerunning simulations if the processed dataset already exists. A new visualization can be created from stored cell-level records. This modularity is one of the reasons the project is suitable as a GitHub research platform rather than only a one-time thesis script.

\section{Research Hypotheses}
Based on the biological background and computational design, several working hypotheses guided the analysis. First, early wall attachment should be associated with later biofilm development because attached cells can seed microcolonies. Second, local signal accumulation and QS activation should influence EPS-producing states and therefore structural maturation. Third, flow and shear should modulate attachment stability and spatial distribution. Fourth, early cell and cluster features should improve prediction compared with simulation parameters alone because they capture the realized trajectory of each stochastic run.

These hypotheses are intentionally phrased as expectations rather than absolute claims. The model may reveal exceptions. For example, high adhesion may not produce a thick biofilm if cell division is weak. High flow may inhibit attachment but still increase transport. EPS fraction may not always correlate with thickness if EPS-producing cells remain sparse. Such exceptions are scientifically useful because they show that the system is not governed by one parameter alone.

\chapter{Methodology}

\section{Biofilm Simulation Framework}
The simulation framework was designed to describe biofilm formation in a near-wall catheter-like environment. The implementation combines an individual cell layer, a flow-field layer, and a chemical signal layer. Although the model is still simplified compared with a full experimental catheter system, it deliberately keeps several biological processes that are central to \textit{Staphylococcus aureus} biofilm development: attachment, growth, division, quorum-sensing activation, EPS production, cluster formation, and flow-related detachment. The overall aim is not to reproduce every molecular detail, but to generate controlled, structured, and analyzable spatiotemporal data.

\subsection{Cellular Potts Model for Cell-Level Behavior}
The Cellular Potts Model (CPM) is used as the conceptual basis for representing cell morphology and local interactions. In a CPM, cells occupy connected regions on a lattice. Their movement and shape changes are governed by an energy function, often called a Hamiltonian. A general form is
\begin{equation}
H = \sum_{i,j} J_{\tau(\sigma_i),\tau(\sigma_j)}(1-\delta_{\sigma_i,\sigma_j}) + \sum_{c}\lambda_V(V_c-V_c^{target})^2,
\end{equation}
where \(\sigma_i\) denotes the cell identity at lattice site \(i\), \(\tau\) is the cell type, \(J\) describes adhesion energy between neighboring types, \(V_c\) is the current area or volume of cell \(c\), and \(V_c^{target}\) is the preferred cell size. Lower effective energy favors configurations where adhesion and volume constraints are more satisfied.

In the implemented project, the original browser simulation in \texttt{src/main.js} preserves this CPM-oriented logic through the Artistoo environment. The Python simulation runner in \texttt{simulation/run\_simulation.py} is an export-oriented companion. It does not replace the original simulation; rather, it provides a deterministic, batch-friendly version that can run parameter sweeps and produce standardized database tables. One possible explanation for this two-layer design is practical: browser-based CPM simulation is useful for visual inspection, while the Python runner is more suitable for generating many complete simulation records for machine learning.

Cell states are represented in a biologically interpretable way. The main exported states are \texttt{planktonic}, \texttt{attached}, \texttt{active}, \texttt{inactive}, \texttt{qs\_active}, \texttt{eps\_producing}, and \texttt{detached}. These states are not merely colors in a plot. They determine which cells contribute to cluster formation, which cells produce higher local signaling levels, and which cells contribute to EPS-related summary variables.

\subsection{Fluid Dynamics and Navier--Stokes Representation}
Fluid flow was included because catheter-associated biofilms are strongly shaped by hydrodynamic conditions. At the continuum level, incompressible Newtonian flow can be described by the Navier--Stokes equations:
\begin{align}
\nabla \cdot \mathbf{u} &= 0,\\
\rho\left(\frac{\partial \mathbf{u}}{\partial t}+\mathbf{u}\cdot\nabla\mathbf{u}\right) &= -\nabla p + \mu\nabla^2\mathbf{u}+\mathbf{f},
\end{align}
where \(\mathbf{u}\) is velocity, \(p\) is pressure, \(\rho\) is density, \(\mu\) is dynamic viscosity, and \(\mathbf{f}\) represents body forces. In the current simulation, the flow field is simplified into a near-wall Poiseuille-like profile, which is sufficient for representing the qualitative relationship between wall distance, velocity, and shear. The model assumes that velocity is lower near the wall and higher away from it. This is a reasonable first approximation for studying initial surface colonization.

The project also contains a lightweight Navier--Stokes server path in \texttt{src/ns\_simple.py}, used by the browser simulation. The flow field is drawn as an overlay and can be exported into the database. In the batch Python runner, the flow field is stored in the \texttt{flow\_field} table as grid-level records containing \(x\), \(y\), \(v_x\), \(v_y\), and shear. For downstream machine learning, these high-dimensional grid records are summarized into statistics such as mean velocity, maximum velocity, mean shear, maximum shear, and low/high shear fractions.

Hydrodynamic studies of biofilm systems suggest that flow can affect both biofilm morphology and detachment \cite{dynamicflow2022,stoodley1994}. It is worth noting that flow does not have a single-direction effect. Stronger flow may inhibit attachment through shear, but it may also improve transport of nutrients or signaling molecules. The simulation therefore allows flow rate to influence attachment and movement, while the later machine learning analysis examines whether the resulting outcomes show consistent patterns.

\subsection{Diffusion, Convection, and Local Signaling}
Besides mechanical forcing, chemical transport is another important component of biofilm evolution. Nutrients, metabolites, and quorum-sensing molecules are not uniformly distributed in a real biofilm. In a simplified form, concentration \(C\) can be modeled as
\begin{equation}
\frac{\partial C}{\partial t} + \mathbf{u}\cdot\nabla C = D\nabla^2 C + S - kC,
\end{equation}
where \(D\) is the diffusion coefficient, \(S\) is local production, and \(k\) is the decay rate. The term \(\mathbf{u}\cdot\nabla C\) represents convection by fluid flow. In the implementation, the signal field is updated through a diffusion and decay process. Cells produce signal locally, and the field spreads over time. The parameter \texttt{diffusion\_rate} controls spatial spreading, while \texttt{signal\_decay} prevents unbounded accumulation.

This design was chosen because an uncontrolled signal field would quickly become biologically uninformative. If diffusion is too high or decay is too weak, the entire domain becomes saturated, and all cells may appear QS-active. Such a result would not help mechanism analysis. By including both diffusion and decay, the model can preserve local signal heterogeneity, which is then reflected in features such as \texttt{mean\_local\_signal}, \texttt{max\_local\_signal}, and \texttt{qs\_active\_ratio}.

\subsection{Attachment, Growth, Division, EPS, and QS}
Initial attachment is modeled as a probabilistic near-wall event. A planktonic cell can attach when its distance to the wall is less than a defined attachment distance. The attachment probability is
\begin{equation}
P_{attach}=a_w \exp\left(-\frac{\tau_{local}}{s}\right),
\end{equation}
where \(a_w\) corresponds to \texttt{adhesion\_wall}, \(\tau_{local}\) is local shear, and \(s\) is \texttt{shear\_scale}. This expression captures an intuitive relation: stronger wall adhesion makes attachment easier, while stronger shear makes attachment less stable. It does not claim to be a complete molecular adhesion model, but it provides a transparent mechanism linking simulation parameters to cell-state transitions.

Cells continue to grow and divide after attachment. This point is important. Attached cells are not treated as frozen particles. They can still increase volume, divide, activate QS, and produce EPS. Division preserves state information so that daughter cells remain consistent with the local biological context. Cell identifiers are generated uniquely, and cluster identifiers are recalculated after each exported timepoint. These bookkeeping details may seem technical, but they determine whether the database can be trusted for later analysis.

EPS production is linked to QS-active or EPS-producing cells. This prevents the EPS fraction from becoming a purely arbitrary output. In the current model, cells with high local signal may enter a QS-active state, and some become EPS-producing. The exported \texttt{eps\_fraction} therefore measures the proportion of cells in an EPS-producing state, while EPS-related spatial effects are represented through cell states and cluster morphology. It is worth noting that this is still a simplified representation of EPS. Real EPS includes polysaccharides, proteins, extracellular DNA, and other components, and future versions could separate these components explicitly.

\section{Simulation Data Generation}
The data used in this thesis were generated through repeated simulation runs. Each run corresponds to a unique \texttt{simulation\_id}. The major input parameters include \texttt{flow\_rate}, \texttt{adhesion\_wall}, \texttt{adhesion\_cell}, \texttt{diffusion\_rate}, \texttt{signal\_decay}, \texttt{qs\_threshold}, \texttt{division\_rate}, \texttt{eps\_rate}, \texttt{runtime}, and \texttt{random\_seed}. These parameters represent the physical, biological, and stochastic conditions of a virtual experiment.

\subsection{Parameter Sweep}
A parameter sweep was used to explore many combinations of input parameters. The grid in \texttt{config.yaml} includes multiple values for flow rate, wall adhesion, cell-cell adhesion, diffusion rate, signal decay, QS threshold, division rate, EPS rate, and random seed. For each combination, a new simulation is executed and exported. This design is close to a controlled factorial experiment: the parameter values are not changed randomly without record, but stored together with each simulation outcome.

Parameter sweep is especially useful in this project because it creates variability. A machine learning model cannot learn a meaningful relationship if all simulations are generated under nearly identical conditions. By varying flow rate, adhesion, diffusion, QS threshold, and EPS rate, the dataset contains different possible growth trajectories. Some simulations produce thicker biofilms, some form fewer clusters, and some reach stable structure earlier.

\subsection{Random Seeds and Repeated Experiments}
Random seeds are included because biofilm formation contains stochastic elements. Initial cell positions, movement noise, division events, and attachment events can vary even when the nominal parameter set is unchanged. Running multiple seeds for the same parameter combination provides a way to estimate how much variation comes from stochastic realization rather than parameter difference. This observation is important for interpretation. If two simulations differ only because of random seed, then a model should not overstate that difference as a deterministic mechanism.

\subsection{Timepoints and Export Interval}
Each simulation evolves over multiple timepoints. Instead of exporting every internal step, the project uses an \texttt{export\_interval}. For example, if the runtime is 100 and export interval is 10, the model records approximately 11 snapshots: timepoints 0, 10, 20, and so on until the final state. At each exported timepoint, the cell table may contain hundreds of rows, one for each cell. The cluster table may contain several rows, one for each detected cluster. The final summary table contains one row for the final state of the simulation.

The size of the generated database follows naturally from this structure. If there are 50 parameter combinations, 5 random seeds, 100 exported timepoints, and 500 cells per timepoint, the cell table contains:
\begin{equation}
50 \times 5 \times 100 \times 500 = 12,500,000
\end{equation}
cell-level records. Therefore, large data volume is not the result of manual entry. It is produced by the multiplication of parameter combinations, stochastic repetitions, timepoints, and cell counts.

\section{Database Design}
A relational database was used to store the simulation output. The database can be managed through DBeaver and can run on either MySQL or PostgreSQL. SQLAlchemy is used in the Python code so that the database layer remains relatively independent of the specific database engine.

\subsection{The \texttt{simulations} Table}
The \texttt{simulations} table stores one row for each simulation run. It contains the unique \texttt{simulation\_id}, all input parameters, runtime, random seed, and creation time. This table is the starting point for most analysis because it defines the experimental condition.

\subsection{The \texttt{cells} Table}
The \texttt{cells} table stores cell-level time-series records. Each row contains \texttt{cell\_id}, \texttt{simulation\_id}, \texttt{timepoint}, \texttt{x}, \texttt{y}, \texttt{state}, \texttt{volume}, \texttt{local\_signal}, and \texttt{cluster\_id}. This is the largest table because every exported timepoint may contain many cells. The table is essential for reconstructing early biofilm dynamics, but it is not used directly as the machine learning sample table. Treating every cell as an independent sample would be biologically misleading because cells from the same simulation are strongly correlated.

\subsection{The \texttt{clusters} Table}
The \texttt{clusters} table records cluster-level structure at each exported timepoint. It contains \texttt{cluster\_id}, \texttt{simulation\_id}, \texttt{timepoint}, \texttt{cluster\_size}, \texttt{center\_x}, \texttt{center\_y}, and \texttt{density}. These variables provide a bridge between individual cell behavior and whole-biofilm morphology. For example, two simulations may have similar cell counts but different cluster organization, which can affect roughness or stability.

\subsection{The \texttt{flow\_field} Table}
The \texttt{flow\_field} table stores spatial grid records for each simulation. Each row includes \texttt{x}, \texttt{y}, velocity components, and shear. This table allows downstream analysis to include fluid effects without recomputing the flow field every time. In feature engineering, the grid is summarized into flow-level descriptors such as mean speed, maximum speed, mean shear, maximum shear, and fractions of low-shear or high-shear grid points.

\subsection{The \texttt{simulation\_summary} Table}
The \texttt{simulation\_summary} table contains the final outcomes used as prediction targets. These include \texttt{final\_cell\_count}, \texttt{biofilm\_area}, \texttt{biofilm\_thickness}, \texttt{surface\_coverage}, \texttt{mean\_cluster\_size}, \texttt{max\_cluster\_size}, \texttt{cluster\_count}, \texttt{roughness}, \texttt{qs\_active\_ratio}, \texttt{eps\_fraction}, \texttt{time\_to\_first\_cluster}, and \texttt{time\_to\_stable\_biofilm}. In the modeling pipeline, these columns act as labels.

\begin{table}[H]
\centering
\caption{Database size used in the final modeling workflow.}
\label{tab:database-size}
\begin{tabular}{lrrrrr}
\toprule
Table & Rows & Columns & Missing values & Duplicate rows & Simulation IDs\\
\midrule
\texttt{simulations} & 1,000 & 12 & 0 & 0 & 1,000\\
\texttt{cells} & 293,547 & 9 & 237,289 & 0 & 1,000\\
\texttt{clusters} & 12,976 & 7 & 0 & 0 & 1,000\\
\texttt{flow\_field} & 3,200,000 & 6 & 0 & 0 & 1,000\\
\texttt{simulation\_summary} & 1,000 & 13 & 0 & 0 & 1,000\\
\bottomrule
\end{tabular}
\end{table}

The missing values in the cell table mostly correspond to \texttt{cluster\_id} for cells that are not assigned to a cluster at a given timepoint. This is not necessarily a data-quality failure. It reflects the biological and algorithmic fact that planktonic or isolated cells may not belong to a biofilm cluster. During feature engineering, such missing cluster identifiers are handled safely.

\section{Feature Engineering}
Feature engineering is the central link between simulation and prediction. The raw output is too detailed for direct simulation-level prediction: a single run may contain thousands of cell rows, multiple cluster rows, and thousands of flow-grid records. The goal is to convert these records into a fixed-length representation where each \texttt{simulation\_id} corresponds to exactly one training sample.

\subsection{Cell-Level Aggregation}
Cells are first grouped by \texttt{simulation\_id} and \texttt{timepoint}. For each group, the following features are computed: \texttt{cell\_count}, \texttt{total\_biomass}, \texttt{mean\_volume}, \texttt{std\_volume}, \texttt{max\_volume}, \texttt{mean\_local\_signal}, \texttt{max\_local\_signal}, \texttt{spatial\_spread\_x}, \texttt{spatial\_spread\_y}, \texttt{biofilm\_thickness}, \texttt{surface\_coverage}, and \texttt{roughness}. State ratios are also computed, including \texttt{planktonic\_ratio}, \texttt{attached\_ratio}, \texttt{active\_ratio}, \texttt{inactive\_ratio}, \texttt{qs\_active\_ratio}, and \texttt{eps\_producing\_ratio}.

This aggregation step is important because it makes the model answer a biofilm-level question. A single cell may be attached or planktonic, large or small, high-signal or low-signal. However, the final outcome of interest is not the fate of one cell. The target is the structure formed by the community. For this reason, the cell table is transformed into descriptors of the population and spatial organization.

\subsection{Cluster-Level Aggregation}
Cluster records are also grouped by \texttt{simulation\_id} and \texttt{timepoint}. The extracted features include \texttt{cluster\_count}, \texttt{mean\_cluster\_size}, \texttt{max\_cluster\_size}, \texttt{std\_cluster\_size}, \texttt{mean\_density}, \texttt{std\_density}, \texttt{mean\_center\_x}, \texttt{mean\_center\_y}, \texttt{cluster\_spread\_x}, and \texttt{cluster\_spread\_y}. Cluster features describe whether cells are dispersed, concentrated into a few large groups, or organized into multiple small colonies.

One possible explanation for including cluster features is that biofilm thickness alone can be misleading. A tall but narrow cluster and a moderately thick but widely spread film may have different clinical implications. Cluster count, density, and spread help capture this distinction.

\subsection{Flow-Field Features}
The flow field is aggregated by \texttt{simulation\_id}. The project extracts \texttt{mean\_velocity\_x}, \texttt{mean\_velocity\_y}, \texttt{std\_velocity\_x}, \texttt{std\_velocity\_y}, \texttt{mean\_speed}, \texttt{max\_speed}, \texttt{mean\_shear}, \texttt{max\_shear}, \texttt{std\_shear}, \texttt{low\_shear\_fraction}, and \texttt{high\_shear\_fraction}. The low-shear fraction is the fraction of grid points below the 25th percentile of shear for that simulation, and the high-shear fraction is the fraction above the 75th percentile.

These features are intentionally coarse. They do not preserve the full spatial field, but they make the flow data usable in tabular machine learning models. It is worth noting that this choice may limit spatial interpretation. A future CNN or graph-based model could use the full field more directly.

\subsection{Early-Timepoint Pivoting}
After aggregation, early timepoints are selected. By default, the first 10 exported timepoints are used. Each timepoint feature is expanded into a column. For example, \texttt{cell\_count} becomes \texttt{t0\_cell\_count}, \texttt{t1\_cell\_count}, \texttt{t2\_cell\_count}, and so on. The same logic is applied to cluster features.

The final training dataset is formed by merging simulation parameters, early cell features, early cluster features, flow features, and final labels. The resulting dataset used in this thesis contains 1,000 rows and 368 columns. After excluding identifiers and label columns, the feature matrix contains 361 input features.

\begin{figure}[H]
\centering
\safeincludegraphics[width=0.92\textwidth]{figures/biofilm_cell_state_snapshot.png}
\caption{Example exported biofilm snapshot after feature-oriented visualization. The vertical axis is converted to distance from wall, so the wall is located at \(y=0\) and biofilm growth extends upward.}
\label{fig:cell-state-snapshot}
\end{figure}


\section{Implementation Workflow and Reproducibility}
The computational implementation was organized as a platform rather than as a single script. This choice was made because the research question requires several different operations that should remain reproducible: running simulations, exporting results, writing to a database, loading data, constructing features, training models, producing figures, and making predictions. If these steps are mixed into one file, it becomes difficult to know whether a result comes from the simulator, the feature engineering code, or the model training routine. Therefore, the project was divided into modules with clear responsibilities.

The \texttt{simulation} folder contains scripts for running one simulation and for running parameter sweeps. The single-run script is useful for testing parameter effects or debugging attachment behavior. The parameter-sweep script is used to generate enough data for machine learning. The \texttt{src} folder contains the main data science pipeline: database connection, database initialization, database writing, data loading, feature engineering, dataset building, EDA, baseline modeling, MLP training, evaluation, prediction, and visualization. The \texttt{app} folder contains the Streamlit application. This organization makes the project easier to explain in a defense because each folder corresponds to one part of the research workflow.

A typical reproducible workflow begins with database initialization. The schema is created from \texttt{src/schema.sql}, and the database can be inspected in DBeaver. The simulation is then run with database export enabled. Every simulation produces a new \texttt{simulation\_id}, and all rows written to the five standard tables include this identifier. After enough simulations have been generated, the modeling pipeline loads the data either from database or from CSV backup files. Feature engineering converts the raw tables into a processed training dataset, which is saved as \texttt{data/processed/training\_dataset.csv}. The saved dataset makes the modeling stage reproducible, while the original database remains available for tracing the source of each feature.

Reproducibility also depends on random seeds. In simulation, \texttt{random\_seed} controls stochastic events such as initial cell placement, movement noise, attachment events, and division events. In model training, a fixed random seed controls the train-test split and model initialization. It is worth noting that reproducibility does not mean that all random variation is removed. Instead, the variation is documented. This is closer to a real experimental design, where repeated experiments under the same nominal condition may not produce exactly the same result.

\section{Coordinate System and Visualization Convention}
A nontrivial technical issue in the project is the coordinate system. In many computer graphics environments, the canvas origin is located at the upper-left corner, with the \(y\)-coordinate increasing downward. In biological interpretation, however, it is more natural to put the wall at \(y=0\) and let biofilm height grow upward. If this conversion is not handled carefully, figures may appear flipped, and conclusions about attachment or thickness could be confusing.

For this reason, the analysis figures convert raw exported coordinates into distance from wall. If \(y_{raw}\) is the exported canvas coordinate and \(y_{wall}\) is the wall coordinate, the plotted vertical coordinate is
\begin{equation}
y_{plot}=y_{wall}-y_{raw}.
\end{equation}
With this convention, attached cells should appear near the horizontal axis, and higher biofilm regions should appear above it. The wall is explicitly drawn as a line at \(y=0\). Spatial cell-state plots and flow-field plots use an equal aspect ratio so that the geometry is not distorted by figure scaling. Heatmaps use \texttt{origin=lower}, which also places the wall at the lower part of the figure.

This visualization convention matters for interpretation. If the raw canvas coordinate were plotted directly, the apparent direction of growth would be inverted. A reader might mistakenly think that the biofilm grows downward from the open flow region. By converting the coordinate system, the figures match the physical interpretation: cells attach to the catheter wall and grow outward into the fluid domain.

\section{Data Integrity Checks During Export}
Data export was treated as part of the methodology rather than as a minor technical detail. Each exported table must include \texttt{simulation\_id}. The \texttt{cells} and \texttt{clusters} tables must also include \texttt{timepoint}. Without these fields, the downstream pipeline cannot correctly group observations. A missing \texttt{simulation\_id} would make it impossible to merge cell features with final outcomes. A missing \texttt{timepoint} would make early-timepoint modeling impossible.

The database writer supports several modes for existing simulations: replace, append, and skip. In the main workflow, replace mode is useful when re-running a simulation with the same identifier during debugging. For parameter sweeps, automatically generated identifiers reduce collision risk. The writer uses pandas DataFrames and SQLAlchemy, which makes the export layer flexible for MySQL and PostgreSQL. It also preserves CSV export as a backup path. This dual design is practical because DBeaver is convenient for inspection, while CSV files are convenient for versioned examples or quick tests.

\section{Algorithmic Summary of One Simulation Run}
Each simulation run follows a repeated update cycle. At the beginning, initial cells are seeded in the near-wall domain. A flow field is generated and stored. At each timepoint, cells produce local signal, the signal field diffuses and decays, cell states are updated according to local signal and wall distance, attachment is evaluated for near-wall cells, cells move under flow and random noise, growth and division are applied, new planktonic cells may enter, and boundary removal is performed for cells leaving the domain. After these updates, cluster assignments are recalculated.

At every export interval, the model records the current cell table and cluster table. At the end of the simulation, the final summary is computed. Biofilm thickness is calculated from the distance of biofilm-state cells from the wall. Surface coverage is estimated by the number of occupied wall-adjacent \(x\) positions. Roughness is computed as the standard deviation of biofilm heights. These definitions are simplified but consistent, which is important for machine learning. A noisy or inconsistent label definition would make prediction difficult and interpretation unreliable.

\section{Why the Method Uses Simulation Before Machine Learning}
A possible question is why machine learning is applied to simulation data instead of directly to experimental data. In an ideal case, real catheter biofilm experiments would provide large numbers of labeled examples under controlled flow, adhesion, diffusion, and EPS conditions. In practice, such datasets are difficult to obtain. Experimental biofilm systems are expensive, time-consuming, and often hard to label at the cell level. Simulation provides a controlled intermediate step. It allows the researcher to generate many virtual experiments, test whether the pipeline works, and identify which kinds of features are likely to be informative.

This does not mean that simulation replaces experimental validation. Rather, simulation helps build the computational framework and form hypotheses. If the model suggests that early attached ratio is strongly associated with final thickness, this can motivate an experimental design that measures early attachment under different wall coatings. If shear features appear important, microfluidic experiments could vary flow rate and compare biofilm morphology. The platform therefore sits between mechanistic modeling and experimental planning.


\section{Software Platform Architecture}
The implemented software platform was built in Python and JavaScript with a focus on reproducibility and usability. The original browser simulation uses JavaScript and Artistoo-style CPM logic. The Python modules provide database export, batch simulation, data loading, feature engineering, model training, evaluation, prediction, and Streamlit interaction. The database can be inspected with DBeaver, which provides a visual way to verify that simulation data are truly being written to relational tables.

The architecture can be understood as a pipeline of contracts. The simulation contract is that every run must produce five standard tables. The database contract is that every row must preserve \texttt{simulation\_id}, and time-series rows must preserve \texttt{timepoint}. The feature engineering contract is that each simulation becomes one row in the modeling dataset. The prediction contract is that new inputs must be aligned with the saved feature template. These contracts prevent hidden inconsistencies. For example, if prediction columns are not aligned with training columns, a model may silently produce invalid predictions. The saved feature template prevents this problem.

\section{Interactive Application Design}
The Streamlit application was designed to make the workflow accessible to users who may not want to run every command manually. The app includes pages for database connection, data import, EDA, model training, cross-validation, ablation study, uncertainty estimation, new parameter prediction, early-data prediction, and model visualization. Each page corresponds to a step in the research workflow.

The database page allows the user to test the MySQL or PostgreSQL connection, create tables, reset tables, show table counts, and preview table rows. This is useful because database errors are common in applied projects. A user can immediately see whether the problem is connection-related, schema-related, or data-related. The data import page supports CSV and Excel uploads, which is helpful when transferring data between machines or when using small examples for demonstration.

The model training page builds the processed dataset and trains models. The prediction pages are especially important for the thesis narrative. In parameter-only prediction, the platform estimates final outcomes for a new condition even if no early simulation data are available. In early-data prediction, the user uploads early cells and clusters, and the platform extracts features automatically. This second mode better reflects the scientific idea of predicting later biofilm development from early observations.

\section{Why Relational Database Instead of Only CSV}
CSV files are simple, but they become inconvenient when the data are multi-table and large. A complete simulation dataset has simulation parameters, cell-level time series, cluster-level time series, flow-grid records, and final summaries. Keeping these as separate CSV files can work, but checking consistency becomes harder. A relational database makes the relationships explicit through \texttt{simulation\_id}.

Using DBeaver also improves transparency. The user can open the database, inspect row counts, preview rows, and confirm that every simulation has corresponding records. This is valuable for a thesis defense because it shows that the dataset is not just a static file but a structured output of repeated simulation. It also makes it easier to explain why millions of rows appear: they are generated by parameter combinations, seeds, timepoints, cells, clusters, and flow-grid points.

\section{Handling Empty or Sparse Cluster Data}
One practical issue in biofilm simulation is that clusters may not exist at early timepoints. If a simulation starts with dispersed planktonic cells, the cluster table may be empty for the first few exports. The feature engineering code must handle this without crashing. In the platform, missing cluster features are filled with zero when appropriate. This reflects the biological meaning: no detected cluster corresponds to zero cluster count, zero mean cluster size, and zero density.

This design avoids a common data-processing error. If an empty cluster table caused a crash, early-stage simulations would be difficult to analyze. If missing clusters were dropped, the dataset would become biased toward simulations that already formed clusters. By safely encoding absence of clusters as zero-valued cluster features, the model can learn that lack of early clusters is itself informative.

\chapter{Exploratory Data Analysis}

\section{Data Scale and Quality}
The final database used for the current thesis version contains 1,000 simulation runs. This scale is large enough to demonstrate the complete modeling workflow, although it should still be regarded as a simulation-generated dataset rather than an experimental clinical dataset. The largest table is \texttt{flow\_field}, with 3.2 million records. The \texttt{cells} table contains 293,547 cell-level records, and the \texttt{clusters} table contains 12,976 cluster-level records. These numbers indicate that the platform is able to manage multi-table, multi-scale output rather than a small hand-written CSV example.

A data-quality check was performed before model training. The main fields in \texttt{simulations}, \texttt{clusters}, \texttt{flow\_field}, and \texttt{simulation\_summary} had no missing values. The \texttt{cells} table had missing values mainly in \texttt{cluster\_id}, because not every cell belongs to a detected cluster. This observation suggests that missingness is partly structural rather than accidental. In later feature extraction, missing cluster assignments are not treated as measurement errors; instead, they are handled as part of the cell-state and cluster-formation process.

\section{Distribution of Final Outcomes}
The distribution of \texttt{biofilm\_thickness} provides a first view of the outcome variability created by the parameter sweep. A useful prediction task requires the target to vary across simulations. If all simulations produced nearly identical thickness, a model might appear accurate without learning a meaningful mechanism. The observed distribution shows that the simulations produced a range of final thickness values, which supports the use of regression models.

\begin{figure}[H]
\centering
\safeincludegraphics[width=0.72\textwidth]{figures/biofilm_thickness_distribution.png}
\caption{Distribution of final biofilm thickness across the 1,000 simulation runs.}
\label{fig:thickness-distribution}
\end{figure}

Roughness is another important structural indicator. A biofilm can be thick but smooth, or thinner but highly uneven. Roughness therefore provides complementary information about spatial heterogeneity. In the model outputs, roughness was strongly predictable by tree-based models, which suggests that the engineered features captured relevant spatial information. However, the MLP did not outperform tree models on this target, indicating that nonlinear capacity alone is not sufficient; the representation and training regime also matter.

\section{Time-Series Analysis}
Time-series plots were generated from the aggregated cell and cluster tables. The mean \texttt{cell\_count} over time describes the overall growth trend. If the curve increases steadily, it indicates proliferation and/or continued inflow of cells. If it plateaus, this may reflect saturation, detachment, or the maximum-cell constraint in the simulation. The plot is not interpreted as a direct experimental growth curve, but as a diagnostic summary of simulated population dynamics.

\begin{figure}[H]
\centering
\safeincludegraphics[width=0.78\textwidth]{figures/cell_count_over_time.png}
\caption{Mean cell count over exported timepoints. The curve summarizes simulation-level population growth rather than individual-cell trajectories.}
\label{fig:cell-count-time}
\end{figure}

Cluster growth is analyzed through \texttt{cluster\_count} and \texttt{mean\_cluster\_size}. These curves help distinguish between two possible growth patterns. In one pattern, the biofilm forms many small clusters; in another, cells consolidate into fewer but larger structures. It is worth noting that both patterns can produce a similar final cell count. Therefore, cluster-level analysis provides information that cell count alone cannot capture.

\begin{figure}[H]
\centering
\safeincludegraphics[width=0.78\textwidth]{figures/cluster_count_over_time.png}
\caption{Cluster count over time. Increasing cluster count indicates the emergence of separated microcolonies, while later stabilization may indicate consolidation or stable attachment.}
\label{fig:cluster-count-time}
\end{figure}

The attached-cell ratio provides a direct view of the early wall-colonization process. Since wall attachment is one of the main biological events in catheter biofilm formation, this feature is especially relevant. A rising attached ratio suggests that planktonic cells are increasingly converted into surface-associated cells. However, attachment is also affected by local shear and stochastic movement, so the trend should be interpreted as a population-level tendency rather than a deterministic rule.

\begin{figure}[H]
\centering
\safeincludegraphics[width=0.78\textwidth]{figures/attached_ratio_over_time.png}
\caption{Attached cell ratio over time. This feature links cell-state dynamics to wall colonization.}
\label{fig:attached-time}
\end{figure}

QS activation and EPS production were also monitored over time. The \texttt{qs\_active\_ratio} describes the proportion of cells whose local signal exceeds the activation threshold. The \texttt{eps\_producing\_ratio} describes the proportion of cells in the EPS-producing state. These curves are useful because biofilm maturation is not only a matter of more cells; it also involves a change in community behavior. When QS-active or EPS-producing fractions increase, it suggests that local density and signal accumulation are beginning to alter cellular states.

\begin{figure}[H]
\centering
\begin{subfigure}{0.48\textwidth}
\centering
\safeincludegraphics[width=\textwidth]{figures/qs_active_ratio_over_time.png}
\caption{QS-active ratio}
\end{subfigure}
\begin{subfigure}{0.48\textwidth}
\centering
\safeincludegraphics[width=\textwidth]{figures/eps_producing_ratio_over_time.png}
\caption{EPS-producing ratio}
\end{subfigure}
\caption{Time evolution of quorum-sensing activation and EPS-producing states.}
\label{fig:qs-eps-time}
\end{figure}

\section{Spatial Structure Analysis}
Spatial visualization is necessary because biofilm morphology cannot be fully understood from scalar features. The cell-state snapshot in Figure~\ref{fig:cell-state-snapshot} shows the location of different cell states relative to the wall. The visualization uses distance from wall on the vertical axis so that biofilm height grows upward. This convention is more biologically intuitive than raw canvas coordinates, where the wall corresponds to high pixel \(y\).

The local signal heatmap shows the spatial distribution of signal intensity. If the entire domain were uniformly high, QS activation would lose explanatory value. In the generated heatmap, signal remains spatially structured, which supports the use of local signal features in prediction. One possible explanation is that diffusion and decay together prevent global saturation, allowing local cell density to influence signal accumulation.

\begin{figure}[H]
\centering
\safeincludegraphics[width=0.85\textwidth]{figures/local_signal_heatmap.png}
\caption{Heatmap of local QS signal. The wall is at the lower boundary after coordinate conversion.}
\label{fig:signal-heatmap}
\end{figure}

The flow and shear field plot provides a physical context for attachment and detachment. The simplified flow field shows spatial variation in shear and velocity, and this information is summarized into flow-field features for machine learning. Even though the current flow representation is simplified, including it is valuable because flow is a plausible driver of biofilm structure.

\begin{figure}[H]
\centering
\safeincludegraphics[width=0.85\textwidth]{figures/flow_shear_field.png}
\caption{Flow and shear field used for simulation-level feature extraction.}
\label{fig:flow-field}
\end{figure}

\section{Parameter Sensitivity Analysis}
Parameter sensitivity analysis was performed by plotting input parameters against final outcomes. Flow rate is expected to affect biofilm thickness through two competing mechanisms. On the one hand, stronger flow can increase shear and make attachment less stable. On the other hand, flow can improve transport and move cells through the domain. In the current simulations, the relationship between \texttt{flow\_rate} and thickness is not interpreted as a universal biological law, but as a model-specific response that helps identify how the implemented assumptions behave.

\begin{figure}[H]
\centering
\safeincludegraphics[width=0.70\textwidth]{figures/flow_rate_vs_biofilm_thickness.png}
\caption{Relationship between flow rate and final biofilm thickness.}
\label{fig:flow-thickness}
\end{figure}

Wall adhesion has a more direct conceptual relationship with attachment. Higher \texttt{adhesion\_wall} increases the probability that near-wall cells become attached, especially under lower shear. This may lead to greater surface coverage or earlier cluster formation. However, it is worth noting that stronger attachment does not automatically imply a thicker biofilm in every simulation. If growth, division, signal activation, or EPS production are limiting, strong wall adhesion may mainly affect early colonization rather than final height.

\begin{figure}[H]
\centering
\safeincludegraphics[width=0.70\textwidth]{figures/adhesion_wall_vs_biofilm_thickness.png}
\caption{Relationship between wall adhesion and final biofilm thickness.}
\label{fig:adhesion-thickness}
\end{figure}

These sensitivity plots serve two purposes. First, they are a visual check that the simulation responds to parameter variation. Second, they help interpret model predictions. If a feature importance plot identifies wall adhesion or early attached ratio as influential, the sensitivity plot provides a biological explanation for why that feature may matter.


\section{Correlation Heatmap and Feature Relationships}
The correlation heatmap was generated from numeric columns in the processed training dataset. Non-numeric identifiers were excluded. Because the feature matrix contains hundreds of early-timepoint features, the plotted heatmap was restricted to a readable subset, including simulation parameters and features most correlated with the selected target. This avoids a common problem in high-dimensional EDA: a complete heatmap with hundreds of variables may look impressive but be nearly impossible to interpret.

\begin{figure}[H]
\centering
\safeincludegraphics[width=0.82\textwidth]{figures/correlation_heatmap.png}
\caption{Correlation heatmap of selected numeric features. The plot excludes non-numeric identifiers and emphasizes parameters and features most associated with the modeling target.}
\label{fig:correlation-heatmap}
\end{figure}

The correlation analysis should not be interpreted as causality. A high correlation between early thickness and final thickness is expected because both describe related geometry at different stages. A correlation between attached ratio and thickness may indicate that early wall colonization supports later growth, but it may also reflect shared dependence on wall adhesion and shear. Therefore, correlation is used as an exploratory tool, not as final proof of mechanism.

\section{Cell Volume and Biomass Dynamics}
Mean cell volume over time provides a rough view of growth and division balance. If cells grow but do not divide, mean volume may increase. If division occurs frequently, volume may remain closer to the target range even as cell count increases. In the simulation, volume is not only a descriptive field; it influences division eligibility. Therefore, mean volume and total biomass are useful early features.

\begin{figure}[H]
\centering
\safeincludegraphics[width=0.75\textwidth]{figures/mean_volume_over_time.png}
\caption{Mean cell volume over time. Changes in mean volume reflect the balance between growth and division.}
\label{fig:mean-volume-time}
\end{figure}

Total biomass was added to the feature set because cell count alone does not capture the amount of biological material. Two simulations may have the same number of cells but different volumes. In a biofilm context, biomass is closer to physical obstruction or matrix-associated burden than cell count alone. This suggests that both count-based and mass-based features should be considered.

\section{Biofilm Thickness as a Time-Resolved Feature}
Although final \texttt{biofilm\_thickness} is a target, early thickness is also a feature. This is not data leakage as long as only early timepoints are used to predict the final outcome. In practice, this mimics a situation where a researcher observes an early simulation or experiment and wants to estimate later development. Early thickness can be viewed as an early warning indicator.

\begin{figure}[H]
\centering
\safeincludegraphics[width=0.75\textwidth]{figures/biofilm_thickness_over_time.png}
\caption{Early biofilm thickness over time, aggregated across simulation runs.}
\label{fig:thickness-time}
\end{figure}

The thickness curve is especially useful when combined with attached ratio and cluster count. A rising attached ratio without increasing thickness may indicate surface colonization without vertical development. Increasing cluster count without increasing thickness may indicate lateral spreading rather than upward growth. Increasing thickness together with cluster size suggests maturation of microcolonies into a more three-dimensional biofilm.

\section{Mean Local Signal and QS Interpretation}
Mean local signal over time helps interpret the QS activation curve. If mean signal rises before the QS-active ratio increases, it suggests that cells are approaching the activation threshold. If QS-active ratio remains low despite increasing cell count, one possible explanation is that signal decay or diffusion prevents local accumulation. This kind of interpretation is useful because it connects numerical features to the biological mechanism represented in the simulator.

\begin{figure}[H]
\centering
\safeincludegraphics[width=0.75\textwidth]{figures/mean_local_signal_over_time.png}
\caption{Mean local signal over time. Signal accumulation provides context for QS activation.}
\label{fig:mean-signal-time}
\end{figure}

It is worth noting that QS behavior in the simulation is simplified. Real \textit{S. aureus} quorum sensing involves molecular pathways and regulatory networks that are not explicitly represented here. The simulated local signal should therefore be interpreted as an abstract communication or density-dependent activation variable. Even so, it is useful for studying how local accumulation can create state transitions and influence EPS-producing behavior.

\section{Additional Parameter Sensitivity Observations}
Besides flow rate and wall adhesion, other parameters also shape the simulated outcomes. Diffusion rate controls how quickly local signal spreads. If diffusion is too low, signals remain highly localized; if it is too high, the field may become smoother and less spatially informative. Signal decay controls how long signal persists. Higher decay may delay QS activation, while lower decay may make activation easier. The QS threshold determines how much signal is needed before cells switch state.

Division rate affects population expansion. A higher division rate can increase cell count and biomass, but only if cells survive and remain within the domain. EPS rate affects the fraction of cells entering EPS-producing states and may influence roughness or stability. Adhesion between cells affects cluster formation because cells that remain close are more likely to be assigned to the same cluster. These parameters do not act independently. Their effects may depend on each other, which is why nonlinear models are reasonable.

A useful example is the interaction between wall adhesion and flow. Strong wall adhesion under low flow may produce rapid attachment. Strong wall adhesion under high flow may still help, but the advantage can be partly offset by shear. Similarly, a low QS threshold may produce early activation, but if cells are too dispersed, local signal may not accumulate enough to create large EPS-producing regions. These interactions are difficult to summarize with one plot, which motivates the use of machine learning and feature importance analysis.

\chapter{Machine Learning Models}

\section{Modeling Objective}
The machine learning objective is to predict final biofilm-level outcomes from simulation parameters, early-timepoint cell and cluster features, and flow-field summaries. The target variables include \texttt{biofilm\_thickness}, \texttt{surface\_coverage}, \texttt{roughness}, \texttt{eps\_fraction}, \texttt{cluster\_count}, and \texttt{time\_to\_stable\_biofilm}. Each row in the training dataset corresponds to one simulation run. This design avoids pseudo-replication from treating cell-level rows as independent training cases.

Before training, missing numeric values are filled using median values or zeros when appropriate. Feature templates are saved so that prediction on new parameters uses the same column order as training. This is a practical but important requirement. If the prediction input has a different column order from the training matrix, model outputs may become meaningless even if the code runs without error.

\section{Linear Regression}
Linear Regression models the target as a linear combination of input features:
\begin{equation}
\hat{y}=\beta_0+\sum_{j=1}^{p}\beta_j x_j.
\end{equation}
It is simple and interpretable, but it assumes an approximately linear relation between features and outcome. In this project, linear regression is included as a baseline. If more complex models only slightly outperform it, then the feature-target relation may be close to linear. If tree-based models perform much better, this suggests nonlinear interactions or threshold-like effects.

One limitation is that the feature matrix contains many correlated early-timepoint features. For example, \texttt{t0\_cell\_count}, \texttt{t1\_cell\_count}, and \texttt{t2\_cell\_count} may be strongly related. Ordinary linear regression can become unstable under such multicollinearity. Therefore, it is useful mainly as a reference rather than the final preferred model.

\section{Ridge Regression}
Ridge Regression adds an \(L_2\) penalty to the linear regression objective:
\begin{equation}
\min_{\beta}\sum_i(y_i-\hat{y}_i)^2+\alpha\sum_j\beta_j^2.
\end{equation}
The penalty shrinks coefficients and can improve stability when features are correlated. In the context of this project, ridge regression is a reasonable middle point between interpretability and regularization. It still assumes linear relations, but it reduces the risk that a small number of correlated features dominate the fit.

\section{Random Forest Regressor}
Random Forest is an ensemble of decision trees trained on bootstrap samples and random subsets of features. It is well suited for nonlinear tabular data because each tree can capture threshold-like behavior. For example, attachment may increase sharply when wall adhesion exceeds a certain value and shear remains low. Such patterns are difficult for a simple linear model but natural for tree-based models.

Random Forest also provides feature importance values. These values should not be interpreted as exact biological causality, but they are useful for mechanism analysis. If early attached ratio, mean local signal, or shear features repeatedly appear as important, this suggests that the corresponding biological process may be influential in the simulation system.

\section{Gradient Boosting Regressor}
Gradient Boosting builds trees sequentially, where each new tree attempts to reduce the errors of the previous ensemble. This often gives strong performance on structured tabular datasets. In the current results, Gradient Boosting performs very well for \texttt{biofilm\_thickness} and \texttt{time\_to\_stable\_biofilm}. One possible explanation is that the model can combine weak nonlinear rules across early-timepoint features, such as early cell count, attachment ratio, and roughness.

However, Gradient Boosting can overfit if the dataset is small or if hyperparameters are not controlled. For this reason, the platform includes cross-validation and optional hyperparameter tuning. In the main thesis analysis, the hold-out results are presented together with the recognition that simulation-generated data may be easier for tree models than real experimental data.

\section{Training and Evaluation Pipeline}
All baseline models use the same feature matrix and target vector. The data are split into training and test sets using a fixed random seed for reproducibility. Evaluation metrics include mean absolute error (MAE), root mean squared error (RMSE), and \(R^2\):
\begin{align}
MAE &= \frac{1}{n}\sum_i |y_i-\hat{y}_i|,\\
RMSE &= \sqrt{\frac{1}{n}\sum_i (y_i-\hat{y}_i)^2},\\
R^2 &= 1-\frac{\sum_i(y_i-\hat{y}_i)^2}{\sum_i(y_i-\bar{y})^2}.
\end{align}
MAE is easy to interpret because it is in the same unit as the target. RMSE penalizes larger errors more strongly. \(R^2\) indicates the fraction of variance explained by the model, although it can be misleading when the target distribution is narrow.

\begin{figure}[H]
\centering
\begin{subfigure}{0.48\textwidth}
\centering
\safeincludegraphics[width=\textwidth]{figures/baseline_biofilm_thickness_pred_vs_true.png}
\caption{Predicted vs true}
\end{subfigure}
\begin{subfigure}{0.48\textwidth}
\centering
\safeincludegraphics[width=\textwidth]{figures/baseline_biofilm_thickness_residuals.png}
\caption{Residuals}
\end{subfigure}
\caption{Baseline model diagnostic plots for biofilm thickness. The residual is defined as true value minus predicted value.}
\label{fig:baseline-diagnostics}
\end{figure}

\section{Feature Importance and Permutation Importance}
Feature importance analysis was used to support mechanism interpretation. Tree-based feature importance reflects how much features reduce prediction error inside the trained trees. Permutation importance measures how much model performance changes when a feature is randomly shuffled. The second method is often easier to interpret because it directly asks whether the trained model still performs well when one feature loses its information.

\begin{figure}[H]
\centering
\begin{subfigure}{0.48\textwidth}
\centering
\safeincludegraphics[width=\textwidth]{figures/biofilm_thickness_gradient_boosting_feature_importance.png}
\caption{Tree feature importance}
\end{subfigure}
\begin{subfigure}{0.48\textwidth}
\centering
\safeincludegraphics[width=\textwidth]{figures/biofilm_thickness_gradient_boosting_permutation_importance.png}
\caption{Permutation importance}
\end{subfigure}
\caption{Feature contribution analysis for biofilm thickness prediction.}
\label{fig:importance}
\end{figure}

Feature importance should be interpreted carefully. A high-importance feature may be a direct mechanism, but it may also be correlated with another mechanism. For example, early thickness may correlate with attached ratio and cluster size. The model may select one feature as important even though the underlying biological process involves several coupled variables. This is why feature importance is used together with EDA plots and biological reasoning.


\section{Fairness of Model Comparison}
A fair comparison between models requires that all models use the same feature matrix, the same target column, and the same train-test splitting strategy. In the implementation, the baseline models and the MLP are trained on the same processed dataset. This does not guarantee perfect fairness because the MLP uses a validation split and early stopping, while the baseline models are fitted differently, but it does ensure that model differences are not caused by different input data.

The comparison also uses several metrics instead of a single score. A model with low MAE but high RMSE may perform well on most cases but fail on a few difficult cases. A model with high \(R^2\) may still have errors that matter biologically if the target is measured on a small scale. This is especially relevant for \texttt{surface\_coverage} and \texttt{eps\_fraction}, where absolute values are small.

\section{Cross-Validation and Hyperparameter Tuning}
The platform includes optional K-fold cross-validation and grid-search tuning. Cross-validation is useful because a single train-test split can be lucky or unlucky. By splitting the data several times, the researcher can estimate how stable model performance is across different subsets. Hyperparameter tuning is useful for Ridge, Random Forest, and Gradient Boosting because their behavior depends on parameters such as regularization strength, number of trees, maximum depth, and learning rate.

However, cross-validation must be used with care. If hyperparameter tuning is performed using the same outer folds used to estimate performance, the results may be optimistic. The implementation separates the main hold-out evaluation from optional cross-validation tools, and the thesis reports the main results as a practical demonstration rather than as a final clinical validation. For a future publication-level study, nested cross-validation would be preferable.

\section{Ablation Study}
Ablation analysis was included to answer a practical question: which groups of features contribute most to prediction? The platform compares parameter-only features, parameter plus flow features, parameter plus early cell features, parameter plus cluster features, and all features together. This analysis is useful because high model accuracy alone does not explain whether the model is mainly using the input parameters or the realized early dynamics.

If parameter-only features already perform well, then the parameter sweep may strongly determine the final outcome. If early cell or cluster features improve performance, then the realized early trajectory contains additional information. This observation is biologically plausible. Even under the same parameters, stochastic early attachment and cluster formation may change the final structure. In that case, early observations are valuable because they capture the realized path rather than only the planned experimental condition.

\section{Interpretability Limits}
Feature importance is helpful, but it does not directly identify causality. In a simulation dataset, some features are generated from related processes. For example, early thickness, roughness, attached ratio, and cluster size can all increase together. A model may assign importance to one of these features because it is a convenient predictor, not because it is the only biological cause. This is why the thesis combines feature importance with EDA plots and mechanistic reasoning.

Another limitation is that feature importance can be affected by correlated variables. If two features carry similar information, importance may be split between them or assigned mostly to one. Permutation importance can also underestimate the value of correlated features because shuffling one feature may not harm performance if another correlated feature remains available. Therefore, interpretability results should be treated as evidence for possible mechanisms, not as final proof.


\section{Model Robustness and Overfitting Considerations}
Overfitting is a central concern in machine learning. A model may perform well on training data but poorly on unseen simulations. In this project, overfitting is controlled through train-test splits, regularization in Ridge Regression, ensemble averaging in Random Forest, boosting control in Gradient Boosting, dropout and early stopping in the MLP, and optional cross-validation. These methods do not eliminate overfitting completely, but they make the evaluation more reliable.

The extremely high \(R^2\) values for some baseline targets should be interpreted with caution. They may indicate that the engineered features are genuinely predictive, but they may also indicate that the simulation rules make the target highly deterministic. For example, if final cluster count is strongly determined by earlier cluster features, then high accuracy is expected. This does not invalidate the model, but it affects how the result should be discussed. The goal is to show that the pipeline works and to analyze within-simulation mechanisms, not to claim direct clinical prediction.

\section{Why Multiple Models Are Needed}
Using several models provides a more balanced view. Linear Regression and Ridge Regression test whether simple additive relationships are sufficient. Random Forest and Gradient Boosting test whether nonlinear threshold-like relations improve prediction. The MLP tests whether a neural function approximator can learn from the engineered features. If all models perform similarly, the feature-target relationship may be simple. If tree models strongly outperform linear models, nonlinear interactions are likely important. If MLP underperforms tree models, the dataset may be better suited to tree ensembles or may require different neural architecture.

This multi-model comparison is more informative than presenting only the best model. In thesis writing, it shows that the result was not selected arbitrarily. It also allows negative or weaker results to be discussed scientifically. For example, the MLP result for surface coverage suggests that not all targets benefit from the same deep learning approach. This observation leads naturally to future work on spatial CNN models.

\section{Choice of Evaluation Targets}
The six prediction targets were selected because they represent different aspects of biofilm structure and dynamics. \texttt{biofilm\_thickness} measures vertical growth. \texttt{surface\_coverage} measures lateral colonization along the wall. \texttt{roughness} measures heterogeneity of height. \texttt{eps\_fraction} represents matrix-associated cell state. \texttt{cluster\_count} measures microcolony organization. \texttt{time\_to\_stable\_biofilm} measures temporal development.

Together, these targets provide a broader view than a single thickness measurement. A biofilm can be thick but narrow, smooth but widespread, or thin but highly EPS-active. By predicting several targets, the platform can describe morphology more comprehensively. This is useful for mechanism analysis because different parameters may affect different targets.

\chapter{Deep Learning Model}

\section{Motivation for Using MLP}
A multi-layer perceptron (MLP) was used because biofilm formation is expected to involve nonlinear relationships. Attachment depends on wall adhesion and shear; QS activation depends on local signal accumulation; EPS production depends on cell state; cluster formation depends on spatial proximity and cell-cell adhesion. These processes interact rather than simply add together. A linear model may capture broad trends, but a neural network can, in principle, approximate more complex functions.

The input to the MLP is the same engineered feature matrix used by the baseline models. This choice makes the comparison fair at the feature level. The MLP is not using raw images or raw cell graphs; it is learning from tabular early-timepoint features. This is a practical first step, but it also explains some limitations. If spatial structure is compressed too much during feature engineering, the MLP cannot recover information that has already been lost.

\section{Network Architecture}
The MLP architecture is:
\begin{equation}
\text{input dimension} \rightarrow 128 \rightarrow 64 \rightarrow 32 \rightarrow \text{output dimension}.
\end{equation}
Each hidden layer uses ReLU activation. Batch normalization and dropout are included in the implemented model. Dropout reduces overfitting by randomly masking hidden units during training, while batch normalization stabilizes the scale of hidden activations. The output dimension is one for single-target regression.

\section{Loss Function and Optimizer}
The model is trained using mean squared error (MSE) loss:
\begin{equation}
L = \frac{1}{n}\sum_i(y_i-\hat{y}_i)^2.
\end{equation}
The optimizer is Adam, which adapts learning rates for individual parameters. A learning-rate scheduler reduces the learning rate when validation loss stops improving. Early stopping is used to stop training when the validation loss no longer improves after a defined patience period.

\section{Training Behavior}
The MLP training curve is useful for diagnosing underfitting or overfitting. If training loss decreases but validation loss increases, the model may be memorizing the training set. If both losses remain high, the model may be underfitting or the features may be insufficient. In the current results, the MLP performs well for some targets such as \texttt{cluster\_count}, \texttt{roughness}, and \texttt{time\_to\_stable\_biofilm}, but less well for \texttt{surface\_coverage} and \texttt{eps\_fraction}. This suggests that some targets may require either more data, better target scaling, or a model that preserves spatial/temporal structure more explicitly.

\begin{figure}[H]
\centering
\safeincludegraphics[width=0.78\textwidth]{figures/mlp_biofilm_thickness_loss_curve.png}
\caption{MLP training and validation loss curve for biofilm thickness prediction.}
\label{fig:mlp-loss}
\end{figure}

\begin{figure}[H]
\centering
\safeincludegraphics[width=0.70\textwidth]{figures/mlp_biofilm_thickness_pred_vs_true.png}
\caption{MLP predicted vs true plot for biofilm thickness.}
\label{fig:mlp-pred}
\end{figure}

It is worth noting that deep learning is not automatically superior in this setting. With 1,000 simulation-level samples and hundreds of engineered features, tree-based models often perform more stably. The MLP still has value because it tests whether a general nonlinear function approximator can learn the mapping, but the current results suggest that the best model depends strongly on the target variable.


\section{Input Scaling and Target Behavior}
Before training the MLP, the input features are standardized using a \texttt{StandardScaler}. This is necessary because the feature matrix contains variables on very different scales: flow rate may be between 0.1 and 1.0, cell count may be tens or hundreds, flow-field records may produce small shear values, and timepoint-derived features may have different units. Neural networks are sensitive to scale because gradient updates depend on the magnitude of inputs.

The targets are not all equally easy for the MLP. \texttt{biofilm\_thickness} and \texttt{roughness} have continuous variation and are related to many early structural features. \texttt{cluster\_count} is also predictable because cluster-derived features are included. In contrast, \texttt{surface\_coverage} and \texttt{eps\_fraction} may occupy narrow numeric ranges. For such targets, a small absolute error can lead to poor \(R^2\), and the MLP may benefit from target transformation or weighted loss.

\section{Why Tree Models Outperformed MLP in Several Cases}
The results show that tree-based models are more stable than the MLP for several targets. This is not a failure of deep learning; it is a reminder that deep learning works best when the architecture matches the data. The current MLP receives a tabular feature vector. It does not know that \texttt{t0\_cell\_count}, \texttt{t1\_cell\_count}, and \texttt{t2\_cell\_count} are ordered in time unless that pattern is learned indirectly. It does not know which features are spatially related. Tree ensembles, by contrast, can quickly learn useful thresholds from tabular data.

One possible explanation is that the simulation rules themselves contain threshold-like operations: attachment occurs when distance from wall is below a threshold; QS activation occurs when local signal exceeds a threshold; cluster detection depends on distance; stable biofilm time depends on a stability criterion. Tree-based models are naturally good at representing such conditional logic. The MLP can approximate it, but may need more data or more careful tuning.

\section{Uncertainty and Bootstrap Confidence Intervals}
The platform includes bootstrap confidence intervals for model metrics. Bootstrap resampling repeatedly samples the test predictions with replacement and recalculates metrics. This gives an approximate interval for MAE, RMSE, and \(R^2\). The interval should not be interpreted as a formal clinical confidence interval, because the data are simulation-generated, but it helps communicate metric stability.

For example, the MLP \(R^2\) for \texttt{biofilm\_thickness} was approximately 0.894 with a bootstrap interval from about 0.860 to 0.918. This suggests that the MLP learns a meaningful signal for thickness, even though it is weaker than the best tree-based baseline. For \texttt{surface\_coverage}, the MLP \(R^2\) was negative, indicating that it performed worse than simply predicting the mean under that evaluation. This observation is worth reporting rather than hiding, because it shows that the analysis is not simply selecting favorable results.


\section{Example Use Case: Screening a New Parameter Combination}
One practical use of the platform is to screen a new parameter combination before running a long simulation. The user enters values such as flow rate, wall adhesion, diffusion rate, QS threshold, and EPS rate. The prediction module builds a feature row using these parameters and fills missing early features from the saved training template. The output is a set of predicted final outcomes. This mode is not as informative as using early cell data, but it is fast.

For example, a user may want to know whether increasing wall adhesion while keeping flow rate moderate is likely to increase final thickness. The platform can provide a preliminary prediction. If the predicted thickness and surface coverage are high, the user may decide to run a full simulation or design a more detailed parameter sweep around that region. In this way, machine learning acts as a guide for simulation planning.

\section{Example Use Case: Predicting From Early Simulation Data}
A more scientifically meaningful use case is early-data prediction. Suppose a simulation has run only through the first 10 exported timepoints. The user uploads early cell and cluster CSV files. The platform aggregates them into early features and predicts final outcomes. This resembles a real experimental scenario: early microscopy or sensor data are available, but the final biofilm state is not yet known.

This mode demonstrates the value of the cell-level database. Because early cells preserve positions, states, volumes, and local signals, the platform can compute attached ratio, QS-active ratio, mean volume, roughness, and cluster features. These features describe how the biofilm is actually developing, not just what parameters were used. If early-data prediction outperforms parameter-only prediction, this suggests that early observations contain meaningful information about future morphology.

\section{How Figures Support the Thesis Argument}
The generated figures are not only decorative. Each one supports a specific part of the argument. Distribution plots show whether targets vary enough for prediction. Time-series plots show whether simulated dynamics are biologically interpretable. Spatial snapshots show whether attached and QS/EPS states occur in reasonable locations. Flow-field plots show how hydrodynamic context is represented. Predicted-vs-true plots show model accuracy, while residual plots reveal systematic errors. Feature importance plots connect prediction to mechanism analysis.

For a final defense, these figures can be organized into a story. First, show the database size to prove that large simulation data were generated. Second, show time-series and spatial plots to demonstrate that the simulation produces interpretable dynamics. Third, show model performance to demonstrate prediction. Fourth, show feature importance and ablation results to support mechanism analysis. This sequence is clearer than presenting model metrics alone.

\chapter{Results}

\section{Database and Dataset Construction Results}
The final database contains 1,000 simulations, 293,547 cell records, 12,976 cluster records, 3,200,000 flow-field records, and 1,000 final summary records. After feature engineering, the processed modeling dataset contains 1,000 rows and 368 columns. The feature matrix contains 361 input features after removing identifiers and target columns. This confirms that the workflow successfully converts a large multi-table simulation database into a fixed-size modeling dataset.

This observation is important for the thesis because it demonstrates the feasibility of the proposed pipeline. The simulation does not only produce pictures; it produces structured data. The database does not only store raw rows; it supports reproducible feature extraction. The machine learning pipeline does not simply read arbitrary CSV columns; it builds biofilm-level features from cell-level and cluster-level observations.

\section{Baseline Model Performance}
Table~\ref{tab:baseline-results} summarizes the best baseline model performance for each target. Tree-based models performed strongly for most targets. Gradient Boosting achieved strong results for \texttt{biofilm\_thickness} and \texttt{time\_to\_stable\_biofilm}. Random Forest performed well for \texttt{surface\_coverage}, \texttt{roughness}, and \texttt{eps\_fraction}. The near-perfect result for \texttt{cluster\_count} should be interpreted carefully because cluster count may be highly determined by engineered cluster features.

\begin{table}[H]
\centering
\caption{Best baseline model performance by target.}
\label{tab:baseline-results}
\begin{tabular}{llrrr}
\toprule
Target & Best model & MAE & RMSE & $R^2$\\
\midrule
\texttt{biofilm\_thickness} & Gradient Boosting & 0.0198 & 0.0252 & 0.9997\\
\texttt{surface\_coverage} & Random Forest & 0.000149 & 0.000311 & 0.9495\\
\texttt{roughness} & Random Forest & 0.00344 & 0.00840 & 0.9997\\
\texttt{eps\_fraction} & Random Forest & 0.00105 & 0.00274 & 0.9658\\
\texttt{cluster\_count} & Linear Regression & $1.13\times10^{-15}$ & $1.98\times10^{-15}$ & 1.0000\\
\texttt{time\_to\_stable\_biofilm} & Gradient Boosting & 0.00717 & 0.0143 & 0.99997\\
\bottomrule
\end{tabular}
\end{table}

The high baseline performance suggests that the engineered features contain strong predictive information about final outcomes. One possible explanation is that early-timepoint features already capture the initial trajectory of biofilm development. For example, if early thickness, attached ratio, cluster count, and local signal are known, the final structure may be highly constrained within the simulation model. Another explanation is that simulation-generated data may be more regular than experimental data. This point should be kept in mind when discussing generalization.

\section{MLP Model Performance}
The MLP results are summarized in Table~\ref{tab:mlp-results}. The MLP achieved good performance for \texttt{biofilm\_thickness}, \texttt{roughness}, \texttt{cluster\_count}, and \texttt{time\_to\_stable\_biofilm}, but it was less stable for \texttt{surface\_coverage} and \texttt{eps\_fraction}. The weaker performance for some targets suggests that a tabular MLP may not be the best model for every output.

\begin{table}[H]
\centering
\caption{MLP model performance by target.}
\label{tab:mlp-results}
\begin{tabular}{lrrr}
\toprule
Target & MAE & RMSE & $R^2$\\
\midrule
\texttt{biofilm\_thickness} & 0.3502 & 0.4555 & 0.8937\\
\texttt{surface\_coverage} & 0.00116 & 0.00145 & -0.0944\\
\texttt{roughness} & 0.0937 & 0.1289 & 0.9352\\
\texttt{eps\_fraction} & 0.0116 & 0.0147 & 0.0167\\
\texttt{cluster\_count} & 0.0657 & 0.0767 & 0.9725\\
\texttt{time\_to\_stable\_biofilm} & 0.2101 & 0.3488 & 0.9841\\
\bottomrule
\end{tabular}
\end{table}

While the MLP model captures nonlinear relationships, tree-based models show more stable performance for certain targets. This is not surprising. Tree ensembles are often strong on medium-sized tabular datasets, especially when the target depends on threshold-like interactions. The MLP may need more samples, stronger regularization, target transformation, or a different architecture to improve performance on small-range targets such as surface coverage.

\section{Model Comparison}
The comparison between baseline models and MLP leads to a more nuanced conclusion than simply saying one model is better. For \texttt{biofilm\_thickness}, the baseline model achieved a much lower RMSE than the MLP. For \texttt{roughness}, both models captured the target, but Random Forest was still stronger. For \texttt{cluster\_count}, both models performed well, suggesting that cluster-related engineered features are highly informative. For \texttt{surface\_coverage} and \texttt{eps\_fraction}, the MLP struggled, possibly because these targets have small numeric ranges and may require better scaling or feature representation.

\begin{figure}[H]
\centering
\safeincludegraphics[width=0.75\textwidth]{figures/ablation_biofilm_thickness.png}
\caption{Ablation study for biofilm thickness. Feature groups were added incrementally to compare their contribution to predictive performance.}
\label{fig:ablation}
\end{figure}

The ablation analysis helps clarify the value of different feature groups. Parameter-only features provide a baseline, but adding early cell and cluster features can improve prediction because they represent realized biological dynamics. This suggests that early simulation output contains information beyond the nominal parameter setting. In other words, two simulations with similar parameters can still diverge due to stochastic events, and early features help capture that divergence.

\section{Interpretation of Results}
The strong performance of tree-based models indicates that the simulation outputs contain structured and learnable relationships. However, it would be too strong to claim that these results directly prove a biological mechanism in real \textit{S. aureus} biofilms. The models learn from simulated data, so their conclusions are conditioned on the assumptions built into the simulator. What they can provide is a mechanism-oriented analysis within the simulated system.

For example, if early attached ratio and early thickness repeatedly appear as important, this suggests that early wall colonization strongly constrains later morphology in the model. If shear-related features are important, this supports the idea that hydrodynamics influences stability. If QS or EPS features matter, this indicates that state transitions are contributing to final structure. These interpretations are meaningful, but they should be validated against experimental data in future work.


\section{Detailed Target-by-Target Interpretation}
For \texttt{biofilm\_thickness}, the best baseline model was Gradient Boosting with an RMSE of 0.0252 and \(R^2=0.9997\). This very high performance suggests that early structural features strongly determine final thickness in the simulated system. The MLP also predicted thickness with positive \(R^2\), but its RMSE was much larger. This difference indicates that the tree ensemble better matched the tabular and threshold-like structure of the data.

For \texttt{surface\_coverage}, Random Forest achieved \(R^2=0.9495\), while the MLP produced a negative \(R^2\). Surface coverage is a small-range target, and its value can depend on fine lateral distribution along the wall. A few cells moving along the wall can change coverage without greatly changing cell count or average volume. This may explain why a tabular MLP struggles. A spatial model that sees occupancy along the wall might be more appropriate.

For \texttt{roughness}, Random Forest achieved excellent performance. Roughness is derived from variation in biofilm height, and early roughness, thickness, cluster structure, and spatial spread are informative predictors. The MLP also achieved a good \(R^2\), but not as high as Random Forest. This suggests that the engineered features preserve enough morphology information for prediction, although not all spatial details are retained.

For \texttt{eps\_fraction}, Random Forest again outperformed the MLP. EPS fraction depends on state transitions and may have sparse or threshold-like behavior. A small number of EPS-producing cells can change the fraction, especially when total cell count is not large. This may make the target noisy relative to the feature representation. It also suggests that future work should model EPS as a spatial matrix field rather than only a cell-state fraction.

For \texttt{cluster\_count}, both baseline and MLP models performed strongly. The near-perfect baseline result should be treated with caution because cluster-derived early features may be highly predictive of final cluster count. This is not necessarily a problem, but it means that cluster count is partly a continuation of already-observed structure. In a real early prediction problem, the usefulness of this target would depend on how early the uploaded cell and cluster data are.

For \texttt{time\_to\_stable\_biofilm}, Gradient Boosting achieved very high performance, and the MLP also performed well. Stability time is computed from simulation dynamics, so early growth and coverage features naturally contain information about it. This result suggests that the platform could be useful for estimating how quickly a simulated biofilm reaches a stable state under a new parameter setting.

\section{Practical Meaning of the Prediction Platform}
The trained models can be used in two prediction modes. In the first mode, the user enters only simulation parameters. Missing early features are filled using the saved feature template. This mode is useful for quick screening, but its uncertainty is higher because the model does not know the realized early dynamics. In the second mode, the user uploads early cell and cluster data. The system extracts early features and predicts final outcomes. This mode is closer to the intended scientific use case because it uses actual early observations.

In a defense demonstration, the database page can show that the data are not manually typed. The data analysis page can show distributions, time-series plots, spatial snapshots, and sensitivity figures. The model training page can show that the dataset is built from database tables. The prediction page can then demonstrate how a new parameter combination is translated into predicted biofilm outcomes. This sequence communicates the central logic of the thesis clearly: simulation generates data, data are engineered into features, models learn predictive relationships, and the platform turns the workflow into an interactive tool.

\section{What the Results Do Not Prove}
The results do not prove that a real catheter biofilm will have exactly the predicted thickness under the same parameter values. The parameter values are simulation parameters, not directly measured clinical variables. The results also do not prove that Gradient Boosting is universally better than MLP for biofilm research. They show that, for the current engineered tabular dataset, tree-based models are highly effective. This distinction matters because an honest discussion of limitations makes the thesis more credible.


\section{Reliability of the Data-Driven Mechanism Analysis}
A key question is whether prediction results can support mechanism analysis. The answer is yes, but only with careful framing. The models do not observe molecular mechanisms directly. They observe engineered variables derived from simulation output. Therefore, a model cannot prove that a biological pathway causes a final structure. What it can do is identify which simulated processes carry predictive information. When early attachment features, cluster features, or local signal features repeatedly improve prediction, this suggests that the corresponding simulated processes are important within the model system.

This distinction makes the analysis more scientifically cautious. For example, if \texttt{attached\_ratio} is important for thickness prediction, the correct conclusion is not simply ``attachment causes thickness'' in all real biofilms. A more accurate statement is that, under the implemented simulation rules, early wall-associated cells are strongly associated with later vertical development. This is still meaningful because it links a measurable early state to a later morphological outcome. It also provides a hypothesis that could be tested experimentally.

The same logic applies to flow and shear features. If shear statistics improve prediction, this suggests that hydrodynamic conditions influence the simulated developmental trajectory. However, the exact magnitude of the effect depends on the flow approximation and attachment probability function. Therefore, the thesis treats hydrodynamic conclusions as model-informed interpretations, not as final physical measurements.

\section{Relationship Between Early Features and Final Outcomes}
The early-feature design is one of the most important parts of the platform. Early cell and cluster features are not arbitrary statistical summaries; they represent different aspects of biofilm development. Cell count and biomass describe population growth. Attached ratio describes wall colonization. QS-active ratio describes signal-mediated state transition. EPS-producing ratio describes matrix-related maturation. Cluster count and cluster size describe spatial organization. Roughness and early thickness describe emerging morphology.

Final outcomes are then predicted from these early descriptors. This resembles how a researcher might examine early experimental images and estimate whether a biofilm is likely to become problematic. A small number of cells near the wall may not be dangerous if they remain isolated and inactive. A similar number of cells may be more concerning if they are attached, locally clustered, and QS-active. The engineered features attempt to quantify these differences.

It is worth noting that early features can sometimes be more informative than input parameters. Parameters describe intended conditions, but early features describe realized behavior. In stochastic simulations, two runs with the same parameter values can differ because cells attach in different locations or divide at different times. Early features capture these realized differences. This is one reason parameter-only prediction is useful but incomplete.

\section{Scientific Meaning of High Baseline Performance}
The high performance of baseline models might raise a concern: are the models too good? In simulation studies, very high \(R^2\) values can occur when the outcome is strongly determined by features derived from the same process. This does not necessarily mean the analysis is invalid. It means the problem is internally consistent and the engineered features are informative. However, it also means that claims should be bounded.

For the current thesis, high baseline performance supports the idea that the simulation-to-feature pipeline is coherent. If the models performed poorly on all targets, it might indicate that the exported data, aggregation, or labels were inconsistent. Strong performance shows that the dataset contains learnable structure. The more important scientific question then becomes which features and processes explain that structure.

At the same time, high performance should not be presented as clinical accuracy. A real biological system contains noise, measurement error, strain variation, and environmental complexity that are not present in the same way in simulation. Therefore, the model results are best understood as proof of platform feasibility and mechanism-oriented analysis within the simulated system.

\section{Interpreting Weak MLP Targets Without Hiding Them}
The MLP did not perform equally well across all targets. This is an important result rather than a problem to hide. In many student projects, deep learning is presented as automatically better, but the current results show a more realistic pattern. For tabular simulation features, tree-based models can outperform a neural network. This observation is useful for the thesis because it demonstrates model comparison rather than model promotion.

The weak MLP performance for \texttt{surface\_coverage} and \texttt{eps\_fraction} suggests that these targets may require different representation. Surface coverage depends on lateral occupancy along the wall. A scalar summary of early features may not fully represent where cells are located. EPS fraction depends on discrete state transitions and may be sensitive to local signal thresholds. A model that sees spatial neighborhoods or temporal sequences may be more suitable.

This discussion naturally motivates the future-work chapter. LSTM or GRU models could learn temporal dynamics directly. CNN models could learn spatial occupancy and signal maps. GNN models could learn cell-cell interactions. In this way, the limitations of the current MLP become a guide for the next research step.

\section{How the Platform Can Be Used in Later Research}
The platform can support several future research workflows. First, it can be used for virtual parameter screening. A researcher can run broad parameter sweeps, identify regions that produce high thickness or high EPS fraction, and then focus detailed simulations on those regions. Second, it can be used for early warning prediction. Given early simulation data, the model can estimate whether the final biofilm is likely to become thick, rough, or stable. Third, it can be used for mechanism comparison. By changing the simulation rules and rerunning the same analysis pipeline, the researcher can compare how different biological assumptions affect predictive features.

This last use is especially interesting. For example, one could modify the attachment probability function, change the EPS production rule, or introduce nutrient limitation. If the important features change, that tells us something about the role of the modified mechanism. The platform therefore provides not only predictions but also a framework for computational experimentation.

\section{Connection Between Figures, Tables, and Narrative}
The final thesis should not present figures as isolated outputs. Each figure should answer a question. The database table answers: how much data were generated? The target distribution answers: is there enough outcome variability? The time-series plots answer: do the simulations show interpretable dynamics? The spatial snapshot answers: are cell states positioned in a biologically reasonable way? The flow-field plot answers: how is hydrodynamic context represented? The model-performance table answers: which models predict well? The feature-importance plot answers: which inputs appear most influential?

When these outputs are connected, the thesis narrative becomes stronger. The reader can follow the path from biological motivation to simulation, from simulation to database, from database to features, from features to models, and from models back to mechanism interpretation. This closed loop is the main contribution of the work.

\section{Recommendations for Presenting the Work in the Defense}
In the final defense, the clearest presentation order is not necessarily the same as the thesis chapter order. A recommended sequence is: first introduce catheter-associated biofilm risk; then show the simulation environment; then show DBeaver with the five tables and row counts; then explain the cell-level to biofilm-level aggregation; then show EDA figures; then show model comparison; finally demonstrate the Streamlit prediction page.

This order helps the audience understand why the project is more than a simulation animation. The DBeaver database makes the data-generation scale visible. The feature-engineering explanation prevents misunderstanding about the training samples. The model results show predictive ability. The Streamlit app shows that the workflow is usable by a researcher rather than only by the developer.

A possible defense sentence is: ``The key idea is that each simulation run becomes one experimental sample. The cells table contains detailed observations, but the model predicts biofilm-level outcomes after aggregation. This avoids treating correlated cells as independent samples and keeps the prediction target biologically meaningful.'' This sentence captures the central logic of the platform.

\chapter{Discussion}

\section{Flow Rate and Biofilm Morphology}
Flow rate is one of the most biologically meaningful parameters in the project. In catheter environments, flow influences both mechanical stress and transport. A low flow rate may allow easier attachment because near-wall cells experience weaker shear. However, low flow may also reduce the transport of nutrients or suspended cells. A high flow rate may increase delivery but can also detach weakly attached cells. This dual effect makes flow difficult to interpret with a simple monotonic statement.

The simulation represents this tradeoff through velocity and shear features. Attachment probability decreases with local shear, while cell movement depends partly on flow. The resulting parameter sensitivity plots show that flow rate has an observable relationship with biofilm thickness, but the relationship should be understood as mediated by other variables, especially wall adhesion and early attachment. One possible explanation is that flow rate changes the probability of establishing stable wall-associated cells, and once a stable early population exists, growth and QS/EPS dynamics become more important.

\section{Adhesion and Initial Attachment}
Wall adhesion directly affects whether cells near the catheter surface become attached. The model uses an attachment probability that increases with \texttt{adhesion\_wall} and decreases with local shear. This reflects the intuition that stronger surface interactions promote colonization, while stronger hydrodynamic forcing destabilizes attachment.

The attached ratio over time is useful because it reveals whether wall colonization occurs early or late. Early attachment can act as a seed for later biofilm growth. Once attached cells divide and produce signal, local density increases, and clusters may form. This observation indicates that early adhesion is not only an initial event; it can set the trajectory for the entire simulation. However, strong adhesion alone is not enough. If division rate is low, QS activation is weak, or EPS production is limited, the final structure may still remain small.

\section{QS, EPS, and Structural Maturation}
QS and EPS are central to biofilm maturation. QS activation reflects local signal accumulation, which is related to cell density and diffusion/decay balance. EPS production reflects a transition toward matrix-forming behavior. In the simulation, EPS-producing cells contribute to the \texttt{eps\_fraction} target and influence the interpretation of biofilm maturity.

It is worth noting that the MLP performed poorly for \texttt{eps\_fraction} compared with Random Forest. One possible explanation is that EPS fraction has a small numeric range and may be sensitive to discrete state transitions. Tree models can handle threshold behavior naturally, while an MLP trained on standardized tabular features may require more data or different target scaling. Another possibility is that the current feature set does not fully capture the spatial context of EPS production. If EPS-producing cells are spatially clustered, a CNN or graph model might represent this better than scalar ratios.

\section{Why Biofilm Forms Heterogeneous Structure}
Biofilm heterogeneity emerges from several coupled processes. Cells do not attach uniformly because their positions and local shear differ. Signal does not accumulate uniformly because production depends on local cell density and diffusion/decay. EPS production is state-dependent, so matrix-like behavior appears only in certain regions. Clusters form where attached or biofilm-state cells are close enough to each other. These interactions can create rough, uneven structures even when the parameter setting is simple.

This suggests that biofilm morphology is not only a direct function of one parameter. Instead, morphology is an outcome of feedback. Attachment changes local density; density changes signal; signal changes cell state; state changes EPS production and cluster stability; cluster growth changes thickness and roughness. The feature engineering pipeline was designed to capture these feedback effects indirectly through early-timepoint descriptors.

\section{Why Model Performance Differs Across Targets}
Different targets have different levels of predictability. Biofilm thickness and roughness are closely connected to early spatial structure, so early-timepoint features provide strong predictive information. Cluster count is directly related to engineered cluster features, which explains its high predictability. Time to stable biofilm is also highly predictable in the simulation because stability is derived from population and coverage dynamics.

Surface coverage and EPS fraction are more difficult for the MLP. Surface coverage is a small-range target and may depend on fine spatial distribution along the wall. EPS fraction depends on state transitions that may be sparse or threshold-like. This may be due to the limited representation of spatial features in the current feature set. The model sees aggregated ratios and statistics, but not the full image of where cells are located relative to each other. This limitation points toward future CNN or graph-based models.

\section{Biological Interpretation}
Within the simulated system, the results support the idea that early biofilm structure is informative for final morphology. Early attached ratio, early thickness, local signal, and cluster features can be interpreted as early indicators of whether the biofilm is likely to mature. This aligns with the biological expectation that once a stable surface-associated microcolony forms, later development becomes easier.

However, there are still limitations. The simulation simplifies many biological details of \textit{S. aureus}, including gene regulation, nutrient consumption, immune interactions, and antibiotic response. The model also treats EPS production as a state-level behavior rather than a chemically detailed matrix. Therefore, the biological conclusions should be presented as simulation-supported hypotheses rather than direct experimental proof.

\section{Data and Model Limitations}
The dataset is generated by simulation, not by direct laboratory observation. This gives excellent control over parameters and labels, but it also means that model performance reflects the assumptions of the simulator. If the simulator over-regularizes behavior, machine learning models may achieve very high \(R^2\) values that would not transfer to real experiments.

The current feature set is tabular. It summarizes spatial and temporal information but does not preserve all geometric details. A biofilm snapshot contains local shape, connectivity, and spatial gradients that cannot be fully represented by mean, standard deviation, or ratio features. This limitation is especially relevant for targets such as roughness and surface coverage.

The current MLP is also relatively simple. It is suitable as a first deep learning model, but it does not explicitly model time order or spatial adjacency. Therefore, it should be seen as a baseline deep learning component rather than the final possible architecture for biofilm prediction.


\section{Connection to Catheter-Associated Infection Risk}
Catheter-associated infection risk depends on whether bacteria can attach, survive, and form a persistent structure. The simulation addresses this process in a simplified environment. Wall attachment represents the first step toward colonization. Growth and division represent expansion. QS and EPS represent maturation. Flow and shear represent physical stress. The final summary variables, such as thickness, surface coverage, and time to stable biofilm, can be interpreted as risk-related structural indicators.

A thicker or more stable simulated biofilm may suggest a condition under which removal becomes more difficult. Higher surface coverage may indicate broader colonization of the catheter wall. Higher EPS fraction may indicate stronger matrix-associated protection. However, these interpretations should remain cautious. Real infection risk also depends on host immune response, antibiotic exposure, catheter material, bacterial strain variation, and protein conditioning films on the surface. The present model captures only part of this system.

\section{Simulation as a Mechanism-Analysis Tool}
The value of simulation is not only prediction. It allows mechanisms to be isolated in a way that is difficult experimentally. For example, one can vary wall adhesion while keeping diffusion, QS threshold, and division rate fixed. One can run repeated random seeds to observe stochastic variation. One can test whether early attached ratio remains predictive when flow rate changes. These controlled comparisons are useful for reasoning about mechanism.

At the same time, simulation can be misleading if its assumptions are forgotten. If attachment probability is defined using a particular exponential relation with shear, then any learned relationship between shear and attachment partly reflects that definition. Therefore, model interpretation must always be tied back to the simulation rules. In this thesis, the models are interpreted as learning within-simulator mechanisms, not as discovering universal laws independently of the simulator.

\section{Data-Driven Modeling as a Bridge}
The data-driven part of the project bridges detailed simulation and interpretable outcomes. A simulation can produce thousands or millions of rows, but a researcher still needs summary evidence. Machine learning models compress these high-dimensional records into predictive relationships. Feature importance and ablation analysis then help identify which groups of variables carry information.

This bridge is useful because it forces the modeler to define measurable features. Terms such as ``biofilm maturity'' or ``structural complexity'' are biologically meaningful but vague. In the platform, they must be represented by variables such as cluster count, roughness, EPS fraction, and time to stable biofilm. This may not capture every nuance, but it makes the analysis explicit and reproducible.

\section{Why Preserving Cell-Level Data Still Matters}
Even though the final machine learning dataset has one row per simulation, preserving cell-level data remains important. Cell-level records allow the researcher to reconstruct time-series curves, spatial snapshots, local signal heatmaps, and state ratios. They also allow future feature engineering without rerunning simulations. For example, if a new feature called ``near-wall active cell density'' becomes interesting, it can be computed from the existing cells table.

This is one reason the relational database design is useful. If only the final feature matrix were stored, much biological detail would be lost. By storing raw cell, cluster, and flow records, the platform remains extensible. Future models can use different aggregation rules, different early-timepoint windows, or even cell-graph representations.

\section{Ethical and Practical Considerations}
Because the project is simulation-based, it does not involve patient data or clinical decision-making. Nevertheless, it is important not to overstate the predictive platform as a medical diagnostic tool. The correct interpretation is that the platform supports mechanism exploration and computational hypothesis generation. Before any clinical application, experimental calibration and validation would be required.

From a practical engineering perspective, the platform also demonstrates how database tools such as DBeaver can support biological modeling. Many biofilm simulations generate output files that are difficult to manage. Storing the data in relational tables makes it easier to query, inspect, and connect to analysis scripts. This is a useful direction for computational biology projects that aim to be reproducible and scalable.


\section{Limitations of the Current Biological Abstraction}
The biological abstraction is useful but incomplete. Real \textit{S. aureus} biofilms involve many molecular and environmental factors not represented in the current model. These include specific adhesins, surface proteins, polysaccharide intercellular adhesin, extracellular DNA, nutrient consumption, metabolic heterogeneity, antibiotic exposure, immune response, and interactions with host proteins. The current simulation compresses many of these processes into state transitions and simplified parameters.

This limitation does not make the model useless. A simplified model can still reveal how attachment, flow, signal, and cluster formation interact. But the interpretation must remain proportional to the model. It is better to say that the results suggest possible mechanisms under the current simulation assumptions than to claim direct proof of clinical behavior.

\section{Limitations of the Current Data Scale}
Although 1,000 simulation-level samples are enough to demonstrate the full pipeline, they are still limited for deep learning. A tabular MLP with hundreds of input features can easily overfit if the dataset is not large and diverse enough. The parameter grid is also discrete, which may make the model perform well at known parameter levels but less reliably between or outside those levels.

A larger dataset could include more parameter values, longer runtimes, more random seeds, and different initial conditions. It could also include deliberately challenging cases, such as high flow with high adhesion, low diffusion with low decay, or high EPS rate with sparse initial cells. Such cases would test whether the models generalize beyond the easiest patterns.

\section{Limitations of Engineered Features}
Feature engineering converts complex spatial and temporal data into fixed-length vectors. This makes machine learning convenient, but it inevitably loses information. The exact shape of a biofilm, the topology of clusters, and local gradients of signal or shear are only partially represented. For some targets, such as roughness, the engineered features work well. For other targets, such as surface coverage or EPS fraction, the lost spatial detail may be more important.

This limitation motivates future CNN and GNN models. CNNs could learn directly from spatial grids, while GNNs could learn from cell-cell neighborhood structure. Such models would reduce the need to manually compress spatial information into scalar features.

\section{Limitations of Evaluation}
The evaluation uses train-test splits and optional cross-validation, but it is still based on simulated data. A stronger evaluation would include external validation: either a new simulation dataset generated from a different parameter distribution or experimental measurements. External validation is important because models can exploit regularities in one dataset that may not hold elsewhere.

Another limitation is that some targets are derived from related features. For example, early cluster features and final cluster count are not independent concepts. High performance on such targets should be interpreted as evidence that the model can continue an observed trajectory, not as evidence that it can infer cluster behavior from unrelated parameters.

\chapter{Conclusion}
This thesis developed a complete workflow for simulation of \textit{Staphylococcus aureus} biofilms and data-driven mechanism analysis. The original biological motivation was preserved: catheter-associated biofilm formation is shaped by cell behavior, fluid dynamics, diffusion, attachment, QS activation, and EPS production. On this basis, the project implemented a reproducible pipeline connecting biofilm simulation, database storage, feature engineering, machine learning, deep learning, visualization, and interactive prediction.

The main contribution is the integration of mechanistic simulation with data-driven modeling. The workflow does not stop at generating visual simulation results. It exports standardized tables, stores them in a relational database, aggregates cell-level and cluster-level observations into biofilm-level features, and trains predictive models for final outcomes. The feature engineering design is especially important because it prevents the common mistake of treating individual cell rows as independent training samples. Instead, each \texttt{simulation\_id} becomes one sample, which matches the biofilm-level targets.

The final database contained 1,000 simulation runs and millions of supporting records across the standard tables. The processed dataset contained 1,000 simulation-level samples and 361 input features. Traditional tree-based models achieved strong performance for most prediction targets, while the MLP captured some nonlinear relationships but was less stable for small-range targets such as surface coverage and EPS fraction. This result suggests that deep learning is useful but not automatically superior; model choice must match data scale, feature representation, and target behavior.

Overall, the platform provides a practical research tool for exploring how simulation parameters and early biofilm dynamics influence final morphology. It can support thesis analysis, presentation figures, and future experimental planning. More importantly, it creates a bridge from mechanism modeling to prediction: biofilm simulation produces structured data, data engineering creates interpretable features, and machine learning helps identify which early conditions are most associated with later biofilm development.

\chapter{Future Work}
Several extensions could improve the scientific and predictive value of the platform.

\section{LSTM, GRU, and Transformer Time-Series Models}
The current model pivots early timepoints into tabular columns. This works, but it does not explicitly model temporal order. Recurrent models such as LSTM and GRU, or attention-based Transformer models, could learn time-dependent patterns more naturally. For example, the rate of increase in attached ratio may matter more than its absolute value at one timepoint. A time-series model could capture such dynamics directly.

\section{CNN Models for Spatial Snapshots}
Biofilm morphology is spatial. The current feature set summarizes spatial information using thickness, roughness, spread, and ratios. A CNN could use spatial images or grids of cell states, signal, EPS, and shear as input. This may improve prediction of targets such as surface coverage and roughness, where local geometry matters.

\section{Graph Neural Networks for Single-Cell Interaction}
Cells interact locally. A graph neural network could represent each cell as a node and neighbor relations as edges. Node features could include state, volume, local signal, and distance from wall. Edge features could include distance or adhesion strength. This would allow the model to learn from the single-cell organization without incorrectly treating cells as independent training samples.

\section{SHAP and Local Explainability}
Feature importance and permutation importance provide global interpretation. SHAP analysis could provide local explanations for individual predictions. For example, if a new parameter combination is predicted to produce a thick biofilm, SHAP values could indicate whether the prediction is driven by wall adhesion, early attached ratio, local signal, or cluster size.

\section{More Realistic Fluid Coupling}
The current flow field is simplified. A future version could include stronger two-way coupling between biofilm geometry and the fluid field. As the biofilm grows, it could alter local velocity and shear; in turn, altered shear could influence detachment and growth. This would make the simulation more realistic but also more computationally expensive.

\section{Experimental Validation}
The most important future step is validation with experimental data. Simulation results are useful for controlled mechanism exploration, but experimental microscopy or catheter-flow data would be needed to evaluate biological realism. A possible strategy is to train initially on simulation data and then fine-tune or calibrate using experimental observations.

\begin{thebibliography}{99}
\phantomsection
\addcontentsline{toc}{chapter}{References}

\bibitem{donlan2002} Donlan, R. M., and Costerton, J. W. Biofilms: survival mechanisms of clinically relevant microorganisms. \textit{Clinical Microbiology Reviews}, 2002.

\bibitem{hallstoodley2004} Hall-Stoodley, L., Costerton, J. W., and Stoodley, P. Bacterial biofilms: from the natural environment to infectious diseases. \textit{Nature Reviews Microbiology}, 2004.

\bibitem{costerton1999} Costerton, J. W., Stewart, P. S., and Greenberg, E. P. Bacterial biofilms: a common cause of persistent infections. \textit{Science}, 1999.

\bibitem{flemming2010} Flemming, H.-C., and Wingender, J. The biofilm matrix. \textit{Nature Reviews Microbiology}, 2010.

\bibitem{stoodley1994} Stoodley, P., deBeer, D., and Lewandowski, Z. Liquid flow in biofilm systems. \textit{Applied and Environmental Microbiology}, 1994.

\bibitem{dynamicflow2022} Dynamic changes in biofilm structures under dynamic flow conditions. \textit{Applied and Environmental Microbiology}, 2022. PMCID: PMC9680615.

\bibitem{chun2022} Chun, A. L. M., Mosayyebi, A., Butt, A., Carugo, D., and Salta, M. Early biofilm and streamer formation is mediated by wall shear stress and surface wettability: a multifactorial microfluidic study. \textit{MicrobiologyOpen}, 2022.

\bibitem{drescher2013} Drescher, K., Shen, Y., Bassler, B. L., and Stone, H. A. Biofilm streamers cause catastrophic disruption of flow with consequences for environmental and medical systems. \textit{Proceedings of the National Academy of Sciences}, 2013.

\bibitem{graner1992} Graner, F., and Glazier, J. A. Simulation of biological cell sorting using a two-dimensional extended Potts model. \textit{Physical Review Letters}, 1992.

\bibitem{glazier1993} Glazier, J. A., and Graner, F. Simulation of the differential adhesion driven rearrangement of biological cells. \textit{Physical Review E}, 1993.

\bibitem{artistoo} Artistoo documentation and CPM simulation framework. Available online. Accessed for model implementation reference.

\bibitem{sklearn} Pedregosa, F. et al. Scikit-learn: machine learning in Python. \textit{Journal of Machine Learning Research}, 2011.

\bibitem{pytorch} Paszke, A. et al. PyTorch: an imperative style, high-performance deep learning library. \textit{Advances in Neural Information Processing Systems}, 2019.

\bibitem{sqlalchemy} SQLAlchemy documentation. Relational database toolkit and object relational mapper for Python. Available online.

\bibitem{streamlit} Streamlit documentation. Interactive data application framework for Python. Available online.

\end{thebibliography}

\appendix
\titleformat{\chapter}{\normalfont\bfseries\Large}{Appendix \thechapter}{0.8em}{}
\titlespacing*{\chapter}{0pt}{1.2ex}{1.0ex}

\chapter{Reproducibility Commands}
This appendix records the main terminal commands used to reproduce the computational workflow. It is intentionally concise, because the detailed methodology has already been described in the main chapters. The purpose of keeping these commands in the thesis is to make the analysis easier to audit: a reader can see how the database tables were initialized, how parameter-sweep data were generated, how the modeling dataset was built, and how the interactive platform was launched.

The commands assume that the project is executed from the repository root directory and that the Python virtual environment has already been created and activated. In the actual implementation, database credentials are read from \texttt{config.yaml} or from environment variables. The same workflow can be repeated with either MySQL or PostgreSQL as long as the database settings are consistent with the DBeaver connection.

\begin{small}
\begin{lstlisting}[language=bash]
# Initialize database tables
.venv/bin/python src/database_init.py

# Run parameter sweep into database
.venv/bin/python simulation/run_parameter_sweep.py \
  --config config.yaml \
  --shuffle \
  --max_runs 1000 \
  --runtime 100 \
  --export_interval 10 \
  --export database

# Build dataset and train models from database
.venv/bin/python src/main.py \
  --source database \
  --target biofilm_thickness \
  --early_timepoints 10

# Launch interactive platform
streamlit run app/streamlit_app.py
\end{lstlisting}
\end{small}

For a short smoke test, the parameter sweep can be limited with \texttt{--max\_runs 3} and a smaller runtime. For thesis-scale training data, the sweep should include enough parameter combinations and random seeds so that the final dataset contains a meaningful range of biofilm outcomes.

\chapter{Database Schema}
The database schema is organized around one central identifier: \texttt{simulation\_id}. Each simulation run produces one parameter row, many cell rows, many cluster rows, many flow-field rows, and one final summary row. This structure is important because the project does not treat individual cells as independent training samples. Instead, cell-level and cluster-level records are aggregated into biofilm-level features before model training.

The table below summarizes the function of each standard table. The complete schema is stored in \texttt{src/schema.sql}; this appendix gives a compact description for readers who want to understand the data organization without reading SQL code.

\begingroup
\small
\setlength{\tabcolsep}{5pt}
\begin{longtable}{>{\raggedright\arraybackslash}p{0.30\textwidth}>{\raggedright\arraybackslash}p{0.60\textwidth}}
\caption{Summary of standard project tables.}\\
\toprule
Table & Description\\
\midrule
\endfirsthead
\toprule
Table & Description\\
\midrule
\endhead
\texttt{simulations} & Input parameter table. One row corresponds to one complete simulation run.\\
\texttt{cells} & Cell-level time-series table. Each row records a cell at one exported timepoint.\\
\texttt{clusters} & Cluster-level time-series table. Each row records a cluster at one exported timepoint.\\
\texttt{flow\_field} & Flow-grid table containing velocity and shear information.\\
\texttt{simulation\_summary} & Final outcome table used as prediction labels.\\
\bottomrule
\end{longtable}
\endgroup

In the final database used for this thesis, these tables contained 1,000 simulation runs, 293,547 cell-level records, 12,976 cluster-level records, 3,200,000 flow-field records, and 1,000 final summary records. This confirms that the analysis was based on automatically generated simulation data rather than manually typed examples.

\end{document}
